\documentclass[11pt,twocolumn,a4paper]{article}

% Page layout
\usepackage[margin=2.2cm]{geometry}

% Essential packages
\usepackage{amsmath,amssymb,amsthm,mathtools}
\usepackage{bm}
\usepackage{graphicx}
\usepackage{hyperref}
\usepackage{natbib}
\usepackage{physics}
\usepackage{booktabs}
\usepackage{array}
\usepackage{multirow}
\usepackage{enumitem}
\usepackage{float}
\usepackage{longtable}
\usepackage{tabularx}

\usepackage{caption}
\captionsetup{
  font=small,          % or footnotesize, scriptsize
  labelfont=footnotesize,        % Keep "Figure 1" bold
  textfont=it          % Optional: italicize caption text
}

% Font and typography
\usepackage[utf8]{inputenc}
\usepackage[T1]{fontenc}
\usepackage{lmodern}
\usepackage{microtype}

% Theorem environments
\newtheorem{theorem}{Theorem}[section]
\newtheorem{lemma}[theorem]{Lemma}
\newtheorem{proposition}[theorem]{Proposition}
\newtheorem{corollary}[theorem]{Corollary}
\theoremstyle{definition}
\newtheorem{definition}[theorem]{Definition}
\newtheorem{axiom}{Axiom}
\newtheorem{example}[theorem]{Example}
\newtheorem{protocol}[theorem]{Protocol}
\theoremstyle{remark}
\newtheorem{remark}[theorem]{Remark}

% Hyperref setup
\hypersetup{
    colorlinks=true,
    linkcolor=blue,
    citecolor=blue,
    urlcolor=blue
}

% Custom commands
\newcommand{\kB}{k_{\mathrm{B}}}
\newcommand{\Sk}{S_{\mathrm{k}}}
\newcommand{\St}{S_{\mathrm{t}}}
\newcommand{\Se}{S_{\mathrm{e}}}
\newcommand{\Sspace}{\mathcal{S}}
\newcommand{\Depth}{\mathcal{M}}
\newcommand{\Real}{\mathbb{R}}
\newcommand{\Natural}{\mathbb{N}}
\newcommand{\Integer}{\mathbb{Z}}
\newcommand{\Hilbert}{\mathcal{H}}
\newcommand{\Trie}{\mathcal{T}}
\newcommand{\Nvib}{N_{\mathrm{vib}}}
\newcommand{\Brot}{B_{\mathrm{rot}}}

\begin{document}

\title{\textbf{Categorical Compound Database: \\
Oscillation-Counted Molecular Identification \\
Through Ternary Phase Space Addressing}}

\author{
Kundai Farai Sachikonye\\[0.3em]
AIMe Registry for Artificial Intelligence\\[0.3em]
\texttt{kundai.sachikonye@bitspark.com}
}

\date{\today}

\maketitle

\begin{abstract}
Molecular database search in traditional cheminformatics requires $O(N \times d)$ pairwise fingerprint comparisons per query, where $N$ is the number of database compounds and $d$ is the fingerprint dimensionality. For PubChem-scale databases ($N \approx 10^8$, $d = 1024$), this amounts to $\sim 10^{11}$ operations per query. We demonstrate that this computational burden arises from an unnatural representation: molecules are encoded as arbitrary bit vectors in an extrinsic descriptor space, forcing similarity to be \emph{computed} rather than \emph{read}. We develop an alternative founded on a single axiom---the Bounded Phase Space Law---from which we derive that every stable molecule is a bounded oscillatory system whose identity is fully characterised by three independent entropy coordinates: knowledge entropy $\Sk$ (spectral energy distribution), temporal entropy $\St$ (timescale span), and evolution entropy $\Se$ (harmonic network density). These three coordinates define a compact metric space $\Sspace = [0,1]^3$ in which molecules are intrinsically addressed. The three-dimensionality of this space makes base-3 the natural encoding: each ternary digit (trit) refines one coordinate dimension, and an interleaved ternary string provides a hierarchical address that is simultaneously a position and a path. We construct a ternary trie---the \emph{categorical compound database}---in which search is $O(k)$ trie traversal for $k$-trit resolution, independent of database size $N$. Fuzzy search reduces to prefix truncation: shorter prefixes yield coarser resolution and more matches, with each trit corresponding to one oscillation-counting observation. We validate on 39 compounds from the NIST Computational Chemistry Comparison and Benchmark Database spanning diatomics through polyatomics. All 39 compounds are uniquely resolved at trit depth 12. Five of six chemically defined families (halomethanes, hydrogen halides, homonuclear diatomics, linear triatomics, bent triatomics) show intra-group ternary cohesion exceeding inter-group similarity (cohesion ratios 1.09--4.38), with chemical groupings emerging at depth 3 without any chemical knowledge being encoded. At depth 12, the ternary trie achieves a $3{,}328\times$ speedup over brute-force fingerprint search ($N = 39$, $d = 1024$); at PubChem scale ($N = 10^8$), the projected speedup exceeds $10^9$. The fundamental insight is that molecular search is not computation but phase space addressing: each trit is one oscillation-counting observation, and the address \emph{is} the search.

\medskip
\noindent\textbf{Keywords:} ternary phase space addressing, S-entropy coordinates, oscillation counting, molecular database search, ternary trie, bounded oscillatory systems, fuzzy search, prefix matching, cheminformatics, computational complexity

\end{abstract}

\tableofcontents
\newpage

%==============================================================================
% PART I: MATHEMATICAL FOUNDATIONS
%==============================================================================

\part{Mathematical Foundations}

%==============================================================================
\section{Introduction}
\label{sec:introduction}
%==============================================================================

\subsection{The Molecular Search Problem}

The central computational task in cheminformatics is molecular search: given a query molecule, find all similar compounds in a database. The standard pipeline proceeds as follows \citep{rogers2010, willett1998}:

\begin{enumerate}[nosep]
\item Represent the molecule as a SMILES string or molecular graph.
\item Compute a fixed-length binary fingerprint (e.g., Morgan/ECFP, 1024 bits).
\item For each database compound, compute the Tanimoto coefficient:
\begin{equation}
T(A, B) = \frac{|A \cap B|}{|A \cup B|}
\label{eq:tanimoto}
\end{equation}
where $A$ and $B$ are the bit vectors of query and target.
\item Rank all compounds by $T$ and return those above a threshold.
\end{enumerate}

This procedure requires $O(d)$ bit operations per pairwise comparison and $O(N)$ comparisons per query, yielding $O(N \times d)$ total cost. For the PubChem database \citep{kim2023} with $N \approx 10^8$ compounds and standard fingerprint length $d = 1024$, a single query requires $\sim 10^{11}$ operations. Even with optimised implementations and pre-filtering, the fundamental scaling is linear in $N$.

The problem is not merely one of engineering. It is a problem of \emph{representation}. Binary fingerprints encode molecular structure in an \emph{extrinsic} descriptor space: the bit vector is an arbitrary projection of molecular topology into $\{0,1\}^d$, and similarity must be \emph{computed} because it is not \emph{implicit} in the representation.

The limitations of this approach are well recognised. Fingerprint collisions occur when structurally dissimilar molecules hash to similar bit patterns \citep{riniker2013}. The choice of fingerprint radius, bit length, and counting mode are all design decisions with no principled basis \citep{bajusz2015}. Virtual screening campaigns routinely require threshold tuning on a per-target basis, and the optimal threshold varies across chemical series. These difficulties arise because the fingerprint representation is extrinsic: it was designed for computational convenience, not derived from the physics of molecular identity.

The fundamental question is: does there exist a representation in which molecular search complexity is independent of database size?

\subsection{Bounded Phase Space and Oscillatory Systems}

We answer this question affirmatively by starting from a physical observation: every stable molecule is a bounded oscillatory system. A molecule consists of $N_{\mathrm{atoms}}$ atoms bound by interatomic potentials that confine the nuclear motion to bounded regions of phase space. The vibrational modes---$\Nvib = 3N_{\mathrm{atoms}} - 6$ for nonlinear molecules, $3N_{\mathrm{atoms}} - 5$ for linear molecules \citep{wilson1955}---are the characteristic oscillations of this bounded system. Each mode has a well-defined frequency $\omega_i$, and the set $\{\omega_i\}$ constitutes the molecule's vibrational fingerprint.

The boundedness of molecular phase space has a deep consequence: by the Poincar\'e recurrence theorem \citep{poincare1890}, almost every trajectory in a bounded measure-preserving system returns arbitrarily close to its initial condition. We prove in Section~\ref{sec:bounded_phase_space} that this recurrence, combined with self-consistency requirements, implies that oscillatory dynamics is the \emph{unique valid mode} for bounded systems. Static equilibria, monotonic evolution, and chaotic dynamics all violate fundamental consistency constraints.

This result transforms the molecular search problem. If every molecule is a bounded oscillatory system, then molecular identity \emph{is} oscillatory structure. The set of vibrational frequencies $\{\omega_i\}$ does not merely \emph{describe} the molecule---it \emph{is} the molecule's dynamical identity. Any representation that faithfully encodes oscillatory structure should therefore provide an intrinsic coordinate system for molecular identity.

\subsection{The Categorical Alternative}

We construct such a representation. Instead of computing arbitrary descriptors and comparing them, we \emph{address} each molecule in its natural coordinate space. The coordinate space is three-dimensional because three independent quantities suffice to characterise bounded oscillatory structure: (1) the distribution of energy across modes, (2) the span of timescales, and (3) the density of harmonic relationships between modes. These three quantities define the S-entropy coordinates $(\Sk, \St, \Se) \in [0,1]^3$.

The three-dimensionality of this space makes base-3 the natural encoding. A ternary digit (trit) takes values $\{0, 1, 2\}$, and each trit refines one of the three coordinate dimensions. An interleaved ternary string provides a hierarchical address: the first trit refines $\Sk$, the second refines $\St$, the third refines $\Se$, the fourth refines $\Sk$ again, and so on. This address is stored in a ternary trie---a tree data structure where each node has exactly three children.

Search in this structure is $O(k)$ trie traversal for $k$-trit resolution, regardless of how many compounds $N$ the database contains. Fuzzy search is prefix truncation: a shorter prefix yields a coarser cell containing more compounds. The resolution is adjustable and principled---each additional trit corresponds to one more oscillation-counting observation.

\subsection{Scope and Contributions}

This paper makes the following contributions:

\begin{enumerate}[nosep]
\item We derive from the Bounded Phase Space Law that every stable molecule occupies a unique point in the S-entropy coordinate space $\Sspace = [0,1]^3$ (Sections~\ref{sec:bounded_phase_space}--\ref{sec:s_entropy}).
\item We prove that ternary encoding is the natural addressing scheme for a three-dimensional coordinate space, and that the interleaved ternary string provides distance-preserving hierarchical addressing (Section~\ref{sec:ternary}).
\item We construct the categorical compound database as a ternary trie with $O(k)$ search independent of database size (Section~\ref{sec:architecture}).
\item We demonstrate fuzzy search at adjustable resolution via prefix matching (Section~\ref{sec:fuzzy_search}).
\item We validate that chemical similarity emerges from ternary proximity: 5/6 chemical families show ternary cohesion without chemical knowledge being encoded (Section~\ref{sec:similarity}).
\item We prove $O(k)$ scaling and quantify speedups over brute-force fingerprint search (Section~\ref{sec:complexity}).
\end{enumerate}

The paper is self-contained: all mathematical foundations (bounded phase space, oscillatory necessity, S-entropy coordinates, ternary encoding) are derived from scratch within this paper.

\subsection{Paper Structure}

Part~I (Sections~\ref{sec:introduction}--\ref{sec:ternary}) establishes the mathematical foundations: the Bounded Phase Space Law, Poincar\'e recurrence, oscillatory necessity, S-entropy coordinates, and ternary encoding. Part~II (Sections~\ref{sec:architecture}--\ref{sec:protocol}) constructs the categorical compound database: the trie data structure, fuzzy search mechanism, and encoding protocol. Part~III (Sections~\ref{sec:test_suite}--\ref{sec:complexity}) validates the framework: compound encodings, fuzzy search demonstrations, chemical similarity analysis, structure prediction, and complexity analysis. Part~IV (Sections~\ref{sec:discussion}--\ref{sec:conclusion}) discusses implications and concludes.

%==============================================================================
\section{Bounded Phase Space}
\label{sec:bounded_phase_space}
%==============================================================================

\subsection{The Axiom}

\begin{axiom}[The Bounded Phase Space Law]
\label{ax:bpsl}
All persistent dynamical systems occupy bounded regions of phase space with finite Liouville measure, and these bounded regions admit hierarchical partitioning into distinguishable subregions.
\end{axiom}

A dynamical system is \emph{persistent} if it maintains its identity over time---it does not disperse, decay, or escape to infinity. The axiom states that every such system is confined to a finite region $\Omega$ of phase space $\Gamma = \{(\mathbf{q}, \mathbf{p}) \in \Real^{2n}\}$ satisfying:
\begin{equation}
\mu(\Omega) = \int_\Omega \frac{d^n q \, d^n p}{h^n} < \infty
\label{eq:finite_measure}
\end{equation}
where $\mu$ denotes the Liouville measure normalised by the phase space cell volume $h^n$, with $h$ being Planck's constant \citep{planck1901}.

The axiom is empirically justified for molecular systems. A stable molecule exists in a potential well: the nuclear coordinates are bounded by the dissociation energy, and the momenta are bounded by the total energy. The phase space volume is finite, measurable, and admits natural partitioning by vibrational mode.

\begin{remark}
The axiom excludes transient systems (dissociating molecules, ionising atoms, unbound scattering states) from the framework. This exclusion is appropriate: transient systems do not have well-defined vibrational spectra, and molecular databases catalogue only persistent species. The axiom therefore matches the domain of applicability precisely.
\end{remark}

The hierarchical partitioning requirement is also empirically satisfied. Vibrational modes naturally partition molecular phase space: each normal mode defines a subspace, and the direct product of modal subspaces gives the full vibrational phase space. Within each mode, energy levels provide a further partition. This nesting---molecule $\to$ modes $\to$ energy levels---is the physical basis for the hierarchical structure that the axiom postulates.

\subsection{Poincar\'e Recurrence}

\begin{theorem}[Poincar\'e Recurrence {\citep{poincare1890}}]
\label{thm:poincare}
Let $(\Omega, \mu, \phi_t)$ be a measure-preserving dynamical system with $\mu(\Omega) < \infty$. Then for any measurable set $A \subset \Omega$ with $\mu(A) > 0$, almost every point $x \in A$ returns to $A$ infinitely often under the flow $\phi_t$.
\end{theorem}

\begin{proof}
Let $B = \{x \in A : \phi_t(x) \notin A \text{ for all } t > 0\}$ be the set of non-returning points. Define $B_n = \phi_{-nT}(B)$ for fixed $T > 0$. By measure preservation, $\mu(B_n) = \mu(B)$ for all $n$. The sets $\{B_n\}_{n=0}^\infty$ are pairwise disjoint: if $x \in B_m \cap B_n$ with $m < n$, then $\phi_{mT}(x) \in B$ and $\phi_{nT}(x) \in B \subset A$, so $\phi_{(n-m)T}(\phi_{mT}(x)) \in A$, contradicting $\phi_{mT}(x) \in B$. Since $B_n \subset \Omega$ for all $n$ and $\mu(\Omega) < \infty$:
\begin{equation}
\infty > \mu(\Omega) \geq \mu\left(\bigcup_{n=0}^\infty B_n\right) = \sum_{n=0}^\infty \mu(B_n) = \sum_{n=0}^\infty \mu(B).
\end{equation}
This infinite sum is finite only if $\mu(B) = 0$. Hence almost every point in $A$ returns to $A$. Repeating the argument from any return time gives infinitely many returns.
\end{proof}

\begin{corollary}[Recurrence Times]
\label{cor:recurrence_times}
For a bounded molecular system with phase space volume $\mu(\Omega)$ in Planck units, the Kac recurrence time \citep{kac1947} for a region $A$ of measure $\mu(A)$ is:
\begin{equation}
\langle T_{\mathrm{rec}} \rangle = \frac{\mu(\Omega)}{\mu(A)}.
\end{equation}
This sets the fundamental timescale for the oscillatory dynamics. For a molecular vibrational mode with frequency $\omega_i$, the recurrence time is the vibrational period $T_i = 2\pi/\omega_i$.
\end{corollary}

\subsection{Oscillatory Necessity}

\begin{theorem}[Oscillatory Necessity]
\label{thm:oscillatory}
For self-consistent dynamical systems in bounded phase space, oscillatory dynamics is the unique valid mode. Static, monotonic, and chaotic alternatives each violate fundamental consistency requirements.
\end{theorem}

\begin{proof}
We eliminate all alternatives by contradiction.

\textbf{Static:} If the system is static, $\phi_t(x) = x$ for all $t$. Then no evolution occurs, no measurement can distinguish time $t$ from time $t'$, and the system has no dynamics. A static system is not a dynamical system.

\textbf{Monotonic:} If any phase space coordinate evolves monotonically (e.g., $q_i(t)$ is strictly increasing), then $q_i(t) \to \pm \infty$ as $t \to \infty$, contradicting boundedness of $\Omega$. Alternatively, if $q_i(t)$ converges monotonically to a limit, the system becomes asymptotically static, reducing to the previous case.

\textbf{Chaotic (sensitive dependence):} If the dynamics exhibits sensitive dependence on initial conditions, then two initially close trajectories diverge exponentially: $|\delta x(t)| \sim |\delta x(0)| e^{\lambda t}$ with Lyapunov exponent $\lambda > 0$. In bounded phase space, this divergence must eventually saturate, leading to mixing. But mixing destroys reproducibility: the system's microstate becomes unpredictable, and it cannot serve as a reproducible dynamical identity. More precisely, chaotic mixing implies ergodicity, so time-averages equal phase-space averages, and all initial conditions produce the same macroscopic behaviour---the system loses its distinguishing identity.

\textbf{Oscillatory:} The system evolves, is bounded, and returns (by Poincar\'e recurrence). The only bounded, recurrent, non-static, non-monotonic dynamics is oscillatory. Oscillatory dynamics preserves identity through its characteristic frequencies: different molecules have different frequency sets $\{\omega_i\}$, providing reproducible distinction.

We have eliminated static (no dynamics), monotonic (violates boundedness), and chaotic (destroys identity). Oscillatory dynamics is the unique survivor: it is bounded, recurrent, non-trivial, and identity-preserving. This is the operational mode of every stable molecule.
\end{proof}

\begin{remark}
The elimination of chaotic dynamics deserves further comment. In bounded Hamiltonian systems, chaos can coexist with regular motion (the KAM scenario). The argument above is not that chaos cannot occur in bounded systems, but that it cannot serve as the \emph{identity-defining} dynamics. A molecule's vibrational frequencies are well-defined precisely because the vibrational motion is regular (oscillatory), not chaotic. Anharmonic corrections introduce weak coupling between modes but do not destroy the mode structure. The normal mode frequencies remain well-defined experimental observables, confirming that the relevant dynamics is oscillatory.
\end{remark}

\subsection{Frequency-Energy Identity}

\begin{corollary}[Discrete Mode Spectrum]
\label{cor:discrete}
A bounded oscillatory system possesses a discrete spectrum of characteristic frequencies $\{\omega_1, \omega_2, \ldots, \omega_N\}$.
\end{corollary}

\begin{proof}
By Theorem~\ref{thm:oscillatory}, the dynamics is oscillatory. The Koopman operator $\hat{K}$ \citep{koopman1931} acting on observables of a bounded system has a discrete spectrum when the phase space is compact. The eigenvalues of $\hat{K}$ are $e^{i\omega_j t}$ for a discrete set of frequencies $\{\omega_j\}$. For molecular vibrational modes, the Bohr--Sommerfeld quantisation condition on bounded domains gives $E_n = (n + 1/2)\hbar\omega$ \citep{bohr1913}, establishing the frequency-energy identity $E = \hbar\omega$.
\end{proof}

\begin{remark}
The discrete mode spectrum is not merely a quantum mechanical result. It follows from the compactness of the bounded phase space: any measure-preserving flow on a compact space has a discrete Koopman spectrum. The quantum mechanical derivation via Schr\"odinger's equation and the classical derivation via the Koopman operator agree on the mode frequencies, as they must---they describe the same bounded oscillatory system. The S-entropy encoding therefore applies regardless of whether one adopts a quantum or classical description.
\end{remark}

\begin{corollary}[Spectral Completeness]
\label{cor:spectral_complete}
The set of fundamental vibrational frequencies $\{\omega_1, \ldots, \omega_{\Nvib}\}$ provides a complete characterisation of the molecule's oscillatory identity within the harmonic approximation.
\end{corollary}

\begin{proof}
Within the Born--Oppenheimer approximation \citep{born1927}, the nuclear dynamics on the ground-state potential energy surface determines the vibrational frequencies. These frequencies are the eigenvalues of the mass-weighted Hessian of the potential energy surface at its minimum. The eigenvalues are a complete invariant of the quadratic (harmonic) approximation to the potential, and different potential surfaces generically produce different eigenvalues.
\end{proof}

\subsection{Implications for Molecular Systems}

Every stable molecule is a bounded oscillatory system. A nonlinear molecule with $N_{\mathrm{atoms}}$ atoms possesses $\Nvib = 3N_{\mathrm{atoms}} - 6$ normal vibrational modes, each with a characteristic frequency $\omega_i$ determined by the molecular potential energy surface \citep{wilson1955}. A linear molecule has $\Nvib = 3N_{\mathrm{atoms}} - 5$ modes. A diatomic has $\Nvib = 1$.

The set of fundamental vibrational frequencies $\{\omega_1, \ldots, \omega_{\Nvib}\}$ constitutes the molecule's \emph{oscillatory fingerprint}---the complete characterisation of its bounded oscillatory structure. Two molecules are vibrationally identical if and only if they share the same frequency set (up to measurement precision). This fingerprint is accessible experimentally through infrared absorption and Raman scattering spectroscopy \citep{wilson1955, herzberg1945}.

The transition from molecular structure to oscillatory fingerprint is the key step: it moves from an arbitrary structural representation (molecular graph, connectivity table, SMILES string \citep{weininger1988}) to an intrinsic dynamical characterisation (frequency spectrum). The oscillatory fingerprint is not a computed descriptor---it is the physical identity of the molecule as a bounded dynamical system.

\begin{remark}
The vibrational frequencies used in this work are fundamental frequencies---the $\omega_i$ corresponding to single-quantum transitions $v = 0 \to 1$ for each normal mode. These are directly observable quantities, tabulated for thousands of molecules in databases such as the NIST CCCBDB \citep{johnson2022} and Shimanouchi tables \citep{shimanouchi1972}. The encoding therefore requires no quantum chemical computation---it starts from experimentally measured quantities.
\end{remark}

The number of vibrational modes itself carries structural information: diatomics have 1 mode, triatomics have 3 or 4, and a molecule with $N_{\mathrm{atoms}}$ atoms has up to $3N_{\mathrm{atoms}} - 6$ modes. The mode count constrains the type of S-entropy coordinates that can be computed (Section~\ref{sec:s_entropy}), providing an automatic classification by molecular complexity.

The distinction between the oscillatory fingerprint and the traditional structural fingerprint is fundamental. A Morgan fingerprint \citep{rogers2010} encodes which substructural fragments are present in a molecular graph. This is a combinatorial, topological description: it asks ``what bonds are connected to what?'' The oscillatory fingerprint, by contrast, encodes the dynamics of the bounded system: it asks ``how does this molecule vibrate?'' Two molecules may have the same substructures but vibrate differently (e.g., cis and trans isomers), or vibrate similarly despite different topologies (e.g., isoelectronic molecules like CO and N$_2$). The oscillatory fingerprint captures a physical identity that the structural fingerprint may miss.

%==============================================================================
\section{S-Entropy Coordinate Space}
\label{sec:s_entropy}
%==============================================================================

\subsection{Three Dimensions of Molecular Identity}

Given a molecule's vibrational frequency spectrum $\{\omega_1, \ldots, \omega_N\}$, we define three independent quantities that characterise its oscillatory structure:

\begin{enumerate}
\item \textbf{Knowledge entropy} $\Sk$: \emph{what} the system is---the distribution of vibrational energy across modes. This is the spectral fingerprint.
\item \textbf{Temporal entropy} $\St$: \emph{when} in its cycle---how many timescales the system's oscillatory dynamics spans. This is the temporal hierarchy.
\item \textbf{Evolution entropy} $\Se$: \emph{how connected}---the density of harmonic relationships between modes. This is the topological connectivity of the frequency network.
\end{enumerate}

\begin{theorem}[Dimensional Sufficiency]
\label{thm:dimensional}
Three coordinates $(\Sk, \St, \Se)$ suffice to characterise bounded oscillatory molecular systems up to partition-depth resolution.
\end{theorem}

\begin{proof}[Proof sketch]
The vibrational spectrum $\{\omega_i\}$ is the complete oscillatory characterisation (Corollary~\ref{cor:discrete}). Any summary statistic of this spectrum is a function of three independent aspects:

(i) The \emph{shape} of the frequency distribution (how energy is distributed across modes), captured by $\Sk$.

(ii) The \emph{scale} of the frequency range (how many orders of magnitude the timescales span), captured by $\St$.

(iii) The \emph{connectivity} of the frequency network (how many rational-ratio relationships exist between mode pairs), captured by $\Se$.

These three aspects are algebraically independent: $\Sk$ depends on normalised frequency ratios, $\St$ depends on the extreme frequencies, and $\Se$ depends on pairwise rational proximity. Changing one does not determine the others. Together, they encode the spectral fingerprint (via $\Sk$), the temporal hierarchy (via $\St$), and the harmonic topology (via $\Se$) of the bounded oscillatory system. Any property derivable from the vibrational spectrum---anharmonic couplings, Fermi resonances, rovibrational structure---can be expressed as a function of these three coordinates at sufficient resolution.

To verify independence, consider two counterexamples. (a) H$_2$O and O$_3$ have similar $\Sk$ (0.944 vs 0.983) and identical $\Se$ (both have 3 modes, all pairs harmonic) but very different $\St$ (0.285 vs 0.151), showing $\St$ is independent of $\Sk$ and $\Se$. (b) CO$_2$ and N$_2$O have similar $\Sk$ (0.897 vs 0.887) and similar $\St$ (0.419 vs 0.442) but different $\Se$ (0.667 vs 0.333), showing $\Se$ is independent of $\Sk$ and $\St$.
\end{proof}

\subsection{Knowledge Entropy \texorpdfstring{$\Sk$}{S\_k}}

\begin{definition}[Knowledge Entropy]
\label{def:s_k}
For a molecule with vibrational frequencies $\{\omega_1, \ldots, \omega_N\}$:

\textbf{Polyatomic case} ($N \geq 2$): Define the normalised frequency distribution
\begin{equation}
p_i = \frac{\omega_i}{\sum_{j=1}^{N} \omega_j}, \quad i = 1, \ldots, N.
\label{eq:frequency_distribution}
\end{equation}
The Shannon entropy \citep{shannon1948} of this distribution is
\begin{equation}
H = -\sum_{i=1}^{N} p_i \log_2 p_i.
\label{eq:shannon}
\end{equation}
The knowledge entropy is the normalised Shannon entropy:
\begin{equation}
\Sk = \frac{H}{\log_2 N} \in [0, 1].
\label{eq:s_k_poly}
\end{equation}

\textbf{Diatomic case} ($N = 1$): The single vibrational frequency is normalised against the highest known fundamental:
\begin{equation}
\Sk = \frac{\omega}{\omega_{\mathrm{ref}}} \in [0, 1], \quad \omega_{\mathrm{ref}} = 4401 \text{ cm}^{-1} \text{ (H}_2\text{)}.
\label{eq:s_k_diatomic}
\end{equation}
\end{definition}

\textbf{Physical interpretation.} $\Sk$ measures how evenly the molecule distributes its vibrational energy across modes. When $\Sk = 0$, all energy is concentrated in a single mode (maximally asymmetric distribution). When $\Sk = 1$, all modes have equal frequency (maximally uniform distribution). Molecules with diverse, well-separated modes have high $\Sk$; molecules dominated by one mode have low $\Sk$.

For the diatomic case, the normalisation against H$_2$ places $\Sk$ on the same $[0,1]$ scale: H$_2$ has $\Sk = 1.0$ (the fastest oscillator), while heavier diatomics have lower $\Sk$ in proportion to their reduced fundamental frequency.

\begin{example}
Water (H$_2$O) has modes at 1595, 3657, 3756 cm$^{-1}$. The normalised distribution is $p = (0.177, 0.406, 0.417)$, giving $H = 1.499$ bits and $\Sk = 1.499/\log_2 3 = 0.944$. The modes are relatively uniform, yielding high $\Sk$.
\end{example}

\begin{example}
Hydrogen chloride (HCl) has a single mode at 2886 cm$^{-1}$. Thus $\Sk = 2886/4401 = 0.656$.
\end{example}

\subsection{Temporal Entropy \texorpdfstring{$\St$}{S\_t}}

\begin{definition}[Temporal Entropy]
\label{def:s_t}
\textbf{Polyatomic case} ($N \geq 2$):
\begin{equation}
\St = \frac{\log(\omega_{\max}/\omega_{\min})}{\log(\omega_{\mathrm{ref,max}}/\omega_{\mathrm{ref,min}})} \in [0, 1]
\label{eq:s_t_poly}
\end{equation}
where $\omega_{\mathrm{ref,max}} = 4401$ cm$^{-1}$ and $\omega_{\mathrm{ref,min}} = 218$ cm$^{-1}$, so $\omega_{\mathrm{ref,max}}/\omega_{\mathrm{ref,min}} \approx 20.2$.

\textbf{Diatomic case} ($N = 1$) with rotational constant $\Brot$:
\begin{equation}
\St = \frac{\log(\omega_{\mathrm{vib}}/\Brot)}{\log(\omega_{\mathrm{ref,max}}/B_{\mathrm{rot,min}}^{\mathrm{ref}})}
\label{eq:s_t_diatomic}
\end{equation}
where $B_{\mathrm{rot,min}}^{\mathrm{ref}} = 0.39$ cm$^{-1}$ is a large-moment-of-inertia reference value.
\end{definition}

\textbf{Physical interpretation.} $\St$ measures the temporal span---how many orders of magnitude the molecule's oscillatory timescales cover. A molecule with modes ranging from 200 to 4000 cm$^{-1}$ spans more timescales than one with modes from 1000 to 1500 cm$^{-1}$. For diatomics, the vibrational-to-rotational frequency ratio plays the same role, measuring how many timescales separate the fast vibrational motion from the slow rotational motion.

\begin{example}
CO$_2$ has modes at 667, 1388, 2349 cm$^{-1}$. Thus $\St = \log(2349/667)/\log(4401/218) = \log(3.52)/\log(20.2) = 1.258/3.006 = 0.419$.
\end{example}

\begin{example}
N$_2$ (diatomic) has $\omega = 2330$ cm$^{-1}$ and $\Brot = 1.99$ cm$^{-1}$. Thus $\St = \log(2330/1.99)/\log(4401/0.39) = \log(1171)/\log(11285) = 7.066/9.332 = 0.757$.
\end{example}

\subsection{Evolution Entropy \texorpdfstring{$\Se$}{S\_e}}

\begin{definition}[Harmonic Proximity]
\label{def:harmonic}
Two vibrational modes $\omega_a$, $\omega_b$ are \emph{harmonically proximate} if there exist integers $p$, $q$ with $1 \leq q \leq p \leq 8$ such that
\begin{equation}
\left|\frac{\max(\omega_a, \omega_b)}{\min(\omega_a, \omega_b)} - \frac{p}{q}\right| < \delta
\label{eq:harmonic_proximity}
\end{equation}
where $\delta = 0.05$ is the harmonic tolerance.
\end{definition}

The choice of $p_{\max} = 8$ includes all simple rational ratios up to the octave and beyond (1/1, 2/1, 3/1, 3/2, 4/3, 5/3, 5/4, 7/4, 8/5, etc.). Higher-order rationals ($p, q > 8$) represent increasingly weak harmonic relationships that are less likely to produce observable Fermi resonances or combination bands. The tolerance $\delta = 0.05$ corresponds to roughly a 5\% deviation from exact rational ratio---comparable to the typical anharmonic shift of molecular vibrations from their harmonic values.

\begin{definition}[Evolution Entropy]
\label{def:s_e}
For a molecule with $N$ vibrational modes:
\begin{equation}
N_{\mathrm{pairs}} = \frac{N(N-1)}{2}, \quad
N_{\mathrm{harmonic}} = \left|\{(i,j) : \omega_i, \omega_j \text{ are harmonically proximate}\}\right|
\end{equation}
\begin{equation}
\Se = \frac{N_{\mathrm{harmonic}}}{\max(N_{\mathrm{pairs}}, 1)} \in [0, 1].
\label{eq:s_e}
\end{equation}
For diatomics ($N = 1$), $\Se = 0$ (no mode pairs exist).
\end{definition}

\textbf{Physical interpretation.} $\Se$ measures harmonic network density---how interconnected the molecule's vibrational modes are through rational frequency relationships. High $\Se$ means many harmonic coincidences (Fermi resonances \citep{fermi1931}, combination bands). Low $\Se$ means the modes are harmonically isolated.

\begin{theorem}[Harmonic Density and Molecular Complexity]
\label{thm:harmonic_complexity}
For molecules with $\Nvib \geq 3$, the harmonic density $\Se$ correlates with spectroscopic complexity: higher $\Se$ implies more Fermi resonances, more combination bands, and richer rovibrational structure.
\end{theorem}

\begin{proof}[Proof sketch]
A Fermi resonance occurs when two vibrational states are nearly degenerate, i.e., when $\omega_a \approx 2\omega_b$ or $\omega_a + \omega_b \approx \omega_c$ \citep{fermi1931}. The condition $|\omega_a/\omega_b - p/q| < \delta$ with $p/q = 2/1$ directly tests for Fermi-type resonance. More generally, combination bands arise from transitions involving multiple modes simultaneously, and these are most prominent when mode frequencies satisfy rational relationships. Thus $\Se$, which counts the fraction of mode pairs satisfying rational proximity, directly measures the potential for spectroscopic complexity.
\end{proof}

\begin{example}
Water (H$_2$O) has modes at 1595, 3657, 3756 cm$^{-1}$. The ratios are $3657/1595 = 2.29 \approx 2/1$, $3756/1595 = 2.35 \approx 7/3$, and $3756/3657 = 1.03 \approx 1/1$. All three pairs are harmonically proximate, giving $\Se = 3/3 = 1.0$.
\end{example}

\begin{example}
All diatomics have $\Se = 0$ because there is only one mode and no pairs to evaluate.
\end{example}

\begin{example}
CO$_2$ has modes at 667, 1388, 2349 cm$^{-1}$. The ratios are $1388/667 = 2.08 \approx 2/1$, $2349/667 = 3.52 \approx 7/2$, and $2349/1388 = 1.69 \approx 5/3$. All three ratios are within $\delta = 0.05$ of a rational $p/q$ with $p, q \leq 8$, giving $\Se = 3/3 = 0.667$. Note that $\Se < 1$ despite all pairs being harmonic because the normalisation denominator is $N_{\mathrm{pairs}} = 3$, which equals $N_{\mathrm{harmonic}} = 2$ (checking more carefully: $2.08 - 2/1 = 0.08 > 0.05$, so the first pair actually fails the strict threshold, and we need to check $|2.08 - p/q|$ for all $p/q$ up to 8: $2.08 \approx 2/1$ with deviation 0.08, but also $\approx 25/12$ which exceeds $q_{\max} = 8$. In practice, the computation gives $N_{\mathrm{harmonic}} = 2$ and $\Se = 2/3 = 0.667$.)
\end{example}

\begin{example}
H$_2$S has modes at 1183, 2615, 2626 cm$^{-1}$. The ratios are $2615/1183 = 2.21 \approx 2/1$ (deviation 0.21, fails), $2626/1183 = 2.22$ (fails), and $2626/2615 = 1.004 \approx 1/1$ (passes). Only one of three pairs is harmonic, giving $\Se = 1/3 = 0.333$. The two S--H stretches are nearly degenerate (a common feature of symmetric hydrides), but neither is in rational proximity to the bending mode.
\end{example}

\subsection{The S-Entropy Space}

\begin{definition}[S-Entropy Space]
\label{def:s_space}
The S-entropy space is $\Sspace = [0,1]^3$ equipped with the Euclidean metric:
\begin{equation}
d(\mathbf{S}_1, \mathbf{S}_2) = \sqrt{(\Sk^{(1)} - \Sk^{(2)})^2 + (\St^{(1)} - \St^{(2)})^2 + (\Se^{(1)} - \Se^{(2)})^2}.
\label{eq:euclidean}
\end{equation}
\end{definition}

\begin{proposition}[Properties of S-Entropy Space]
\label{prop:s_space_properties}
The S-entropy space $\Sspace = [0,1]^3$ satisfies:
\begin{enumerate}[nosep]
\item \textbf{Compactness:} $\Sspace$ is closed and bounded in $\Real^3$, hence compact by the Heine--Borel theorem.
\item \textbf{Completeness:} $(\Sspace, d)$ is a complete metric space (closed subset of $\Real^3$).
\item \textbf{Uniqueness:} Every molecule maps to a unique point in $\Sspace$ (up to measurement precision of the vibrational frequencies).
\item \textbf{Continuity:} Nearby points correspond to spectrally similar molecules---small changes in vibrational frequencies produce small changes in $(\Sk, \St, \Se)$.
\end{enumerate}
\end{proposition}

\begin{proof}
Properties (1) and (2) follow from standard topology. For (3), distinct vibrational spectra generically produce distinct $(\Sk, \St, \Se)$ values because $\Sk$, $\St$, $\Se$ depend on different independent aspects of the spectrum (distribution shape, frequency range, harmonic connectivity). For (4), $\Sk$ and $\St$ are continuous functions of the vibrational frequencies (Shannon entropy is continuous in the probability simplex, logarithms are continuous on $(0, \infty)$). The harmonic pair count for $\Se$ changes only at the discrete thresholds where $|\omega_a/\omega_b - p/q| = \delta$, so $\Se$ is piecewise constant with respect to continuous frequency variation, hence piecewise continuous.
\end{proof}

\begin{remark}
The metric $d$ on $\Sspace$ is Euclidean, treating the three coordinates as equally important. One could instead use a weighted metric $d_w(\mathbf{S}_1, \mathbf{S}_2) = \sqrt{w_k(\Delta \Sk)^2 + w_t(\Delta \St)^2 + w_e(\Delta \Se)^2}$ with weights reflecting the relative importance of each coordinate for a particular application. The ternary encoding is agnostic to the metric: it partitions the coordinate space uniformly regardless of the metric used for distance computation. The metric affects only the \emph{interpretation} of proximity, not the \emph{structure} of the trie.
\end{remark}

\begin{figure*}[!htbp]
    \centering
    \includegraphics[width=\textwidth]{figures/panel_1_sentropy_space.png}
    \caption{\textbf{S-Entropy Compound Space.}
    (\textbf{A})~All 39 NIST compounds embedded in three-dimensional S-entropy space $\mathcal{S} = [0,1]^3$ with axes $\Sk$ (knowledge entropy), $\St$ (temporal entropy), $\Se$ (evolution entropy). The unit cube wireframe bounds the space. Point size scales with molecular mass; colour encodes molecular type: diatomic (blue), triatomic (green), tetrahedral (orange), polyatomic (red). Diatomics cluster along the $\Se = 0$ plane (no mode pairs, hence no harmonic edges), while polyatomics occupy the interior with $\Se > 0.3$. The clear separation along $\Se$ demonstrates that harmonic network density is the primary discriminant between molecular complexity classes.
    (\textbf{B})~Distribution of knowledge entropy $\Sk$ across the 39 compounds, displayed as stacked histograms coloured by molecular type. Diatomics (blue) span the full $[0.13, 1.00]$ range because their single-mode frequency varies widely ($\mathrm{Cl}_2$ at 560~cm$^{-1}$ to $\mathrm{H}_2$ at 4401~cm$^{-1}$). Polyatomics (red, orange) concentrate near $\Sk \approx 0.9$--$1.0$, reflecting the high Shannon entropy of multi-mode frequency distributions approaching uniformity.
    (\textbf{C})~Temporal entropy $\St$ versus knowledge entropy $\Sk$ for all compounds. Point size scales with molecular mass; colour encodes type. Diatomics occupy the upper-left quadrant (moderate-to-high $\St$, variable $\Sk$), reflecting their large vibrational-to-rotational frequency ratios. Polyatomics cluster in the lower-right quadrant (high $\Sk$, moderate $\St$), where frequency distributions are uniform but the ratio between highest and lowest modes is moderate. The anti-correlation between $\Sk$ and $\St$ for diatomics arises because light diatomics (high $\omega$, high $\Sk$) tend to have moderate rotational constants (moderate $\St$).
    (\textbf{D})~Evolution entropy $\Se$ for all 39 compounds, sorted in decreasing order. Bar colour encodes molecular type. The distribution is clearly bimodal: diatomics at $\Se = 0$ (leftmost bars) and polyatomics/triatomics at $\Se \geq 0.33$. $\mathrm{H}_2\mathrm{O}$ achieves the maximum $\Se = 1.0$ (all three mode pairs are harmonically proximate). The gap between 0 and 0.33 reflects the structural transition from single-mode to multi-mode molecules.}
    \label{fig:sentropy_space}
\end{figure*}

%==============================================================================
\section{Ternary Encoding}
\label{sec:ternary}
%==============================================================================

\subsection{Why Base 3?}

The S-entropy space is three-dimensional. A natural question is: what is the optimal base for encoding coordinates in this space?

A ternary digit (trit) takes values $\{0, 1, 2\}$, which maps directly to the three coordinate dimensions:
\begin{equation}
\text{trit value} = \begin{cases}
0 & \to \text{refinement in lower third} \\
1 & \to \text{refinement in middle third} \\
2 & \to \text{refinement in upper third}
\end{cases}
\label{eq:trit_values}
\end{equation}
and the position within the interleaved string determines which dimension is being refined:
\begin{equation}
\text{position } j \bmod 3 = \begin{cases}
0 & \to \text{refine } \Sk \\
1 & \to \text{refine } \St \\
2 & \to \text{refine } \Se
\end{cases}
\label{eq:dimension_assignment}
\end{equation}

\begin{theorem}[Ternary Naturalness]
\label{thm:ternary_natural}
For a $d$-dimensional coordinate space $[0,1]^d$, base-$d$ representation provides native $d$-dimensional addressing with each digit refining one coordinate. For $d = 3$ (S-entropy space), base-3 is the natural encoding.
\end{theorem}

\begin{proof}
In base-$d$ representation, each digit takes values $\{0, 1, \ldots, d-1\}$, providing exactly $d$ choices per refinement step. With interleaving across $d$ dimensions (position $j \bmod d$ refines dimension $j \bmod d$), each digit contributes one refinement to one dimension. After $d$ consecutive digits, every dimension has been refined exactly once, producing a cell of volume $d^{-d}$ in the unit hypercube. This is the coarsest refinement that touches all dimensions equally.

In contrast, base-2 (binary) requires $\lceil \log_2 3 \rceil = 2$ bits per trit to encode three choices, wasting $2 - \log_2 3 \approx 0.42$ bits per refinement. Base-4 or higher would require more digits to achieve the same per-dimension resolution. Base-$d$ is optimal in the sense that each digit carries exactly one refinement's worth of information about one dimension.
\end{proof}

\begin{remark}
The choice of base 3 is not arbitrary or merely convenient. It arises from the dimensionality of the S-entropy coordinate space, which in turn arises from the three independent aspects of bounded oscillatory structure. Binary encoding would impose a 1D structure on a 3D space; ternary encoding preserves the native 3D geometry.
\end{remark}

\begin{remark}
Ternary computing has a distinguished history. The Setun computer at Moscow State University \citep{brusentsov1962} demonstrated that balanced ternary arithmetic is more efficient than binary for certain operations. The connection here is different: we do not advocate replacing binary hardware with ternary hardware, but rather observe that the natural representation of a 3D coordinate space is base-3, regardless of the hardware's native base. The ternary trie is implemented on standard binary hardware; the ternary structure is in the \emph{data}, not the \emph{machine}.
\end{remark}

The information-theoretic content of a single trit is $\log_2 3 \approx 1.585$ bits. A binary encoding of the same coordinate space would require 2 bits per refinement step (since $\lceil \log_2 3 \rceil = 2$), wasting $2 - 1.585 = 0.415$ bits per step. Over $k$ refinements, the binary representation wastes $0.415k$ bits---roughly 26\% overhead. More importantly, the binary representation obscures the three-dimensional structure: with 2 bits per step, the mapping from bit position to coordinate dimension requires an additional layer of interpretation that the ternary encoding provides natively.

\subsection{The Interleaving Scheme}

\begin{definition}[Interleaved Ternary Address]
\label{def:interleaving}
Given coordinates $(\Sk, \St, \Se) \in [0,1]^3$, the $k$-trit interleaved ternary address $\mathbf{t} = (t_1, t_2, \ldots, t_k)$ is generated by Algorithm~1.
\end{definition}

\begin{figure}[H]
\begin{center}
\fbox{\parbox{0.85\textwidth}{
\textbf{Algorithm 1: Interleaved Ternary Encoding}
\begin{enumerate}[nosep]
\item[] \textbf{Input:} $(\Sk, \St, \Se) \in [0,1]^3$, depth $k$
\item[] \textbf{Output:} Trit string $\mathbf{t} = (t_1, \ldots, t_k)$
\item[] Initialise: $r_0 \leftarrow \Sk$, $r_1 \leftarrow \St$, $r_2 \leftarrow \Se$
\item[] \textbf{for} $j = 1$ \textbf{to} $k$:
\item[] \quad $\mathrm{dim} \leftarrow (j-1) \bmod 3$
\item[] \quad $t_j \leftarrow \lfloor 3 \cdot r_{\mathrm{dim}} \rfloor$
\item[] \quad $t_j \leftarrow \min(t_j, 2)$ \hfill (clamp to valid trit)
\item[] \quad $r_{\mathrm{dim}} \leftarrow 3 \cdot r_{\mathrm{dim}} - t_j$
\end{enumerate}
}}
\end{center}
\label{alg:encoding}
\end{figure}

Position $j \bmod 3$ determines which dimension is refined. After $3m$ trits, each dimension has been refined $m$ times, producing cells of volume $3^{-3m}$ in the unit cube. Each step extracts one base-3 digit from the remaining fractional part of the corresponding coordinate.

\begin{example}
For H$_2$O with $(\Sk, \St, \Se) = (0.9442, 0.2850, 1.0000)$:
\begin{itemize}[nosep]
\item $j=1$: dim=0 ($\Sk$), $3 \times 0.9442 = 2.833$, $t_1 = 2$, remainder $= 0.833$
\item $j=2$: dim=1 ($\St$), $3 \times 0.2850 = 0.855$, $t_2 = 0$, remainder $= 0.855$
\item $j=3$: dim=2 ($\Se$), $3 \times 1.0000 = 3.000$, $t_3 = 2$ (clamped), remainder $= 0.000$
\end{itemize}
Continuing yields the address \texttt{202222112122\ldots}
\end{example}

\begin{example}
For HCl with $(\Sk, \St, \Se) = (0.6558, 0.6029, 0.0000)$:
\begin{itemize}[nosep]
\item $j=1$: dim=0 ($\Sk$), $3 \times 0.6558 = 1.967$, $t_1 = 1$, remainder $= 0.967$
\item $j=2$: dim=1 ($\St$), $3 \times 0.6029 = 1.809$, $t_2 = 1$, remainder $= 0.809$
\item $j=3$: dim=2 ($\Se$), $3 \times 0.0000 = 0.000$, $t_3 = 0$, remainder $= 0.000$
\end{itemize}
The first three trits are \texttt{110}. All subsequent $\Se$-dimension trits are 0 (since $\Se = 0$ for diatomics). The remaining trits refine $\Sk$ and $\St$ only, producing the full address \texttt{110220210200\ldots}. Note that the zero $\Se$ coordinate immediately marks this as a diatomic in the trie structure.
\end{example}

\subsection{Trit-Cell Bijection}

\begin{definition}[Level-$k$ Cell]
\label{def:cell}
A level-$k$ cell $C_k(\mathbf{t})$ is the region of $\Sspace$ addressed by the $k$-trit string $\mathbf{t} = (t_1, \ldots, t_k)$. It is a rectangular parallelepiped in $[0,1]^3$ with side lengths $3^{-m_0} \times 3^{-m_1} \times 3^{-m_2}$, where $m_i$ is the number of trits assigned to dimension $i$ among $t_1, \ldots, t_k$.
\end{definition}

\begin{theorem}[Trit-Cell Bijection]
\label{thm:bijection}
The mapping $\varphi_k: \{0,1,2\}^k \to \mathcal{C}_k$ from $k$-trit strings to level-$k$ cells is a bijection. Each trit string addresses exactly one cell; each cell has exactly one address.
\end{theorem}

\begin{proof}
By induction on $k$.

\textbf{Base case} ($k = 0$): $\varphi_0(\varepsilon) = \Sspace$ (the whole space). There is one empty string and one cell.

\textbf{Inductive step:} Suppose $\varphi_k$ is a bijection. The cell $C_k(\mathbf{t})$ is subdivided into three equal parts along dimension $(k \bmod 3)$. The trit $t_{k+1} \in \{0, 1, 2\}$ selects one of these three parts to form $C_{k+1}(\mathbf{t}, t_{k+1})$. This is injective because different trit values select different subregions. It is surjective because every subregion has a label in $\{0, 1, 2\}$. Since $\varphi_k$ is a bijection and the extension is a bijection from $\{0,1,2\}$ to the three subregions of each level-$k$ cell, $\varphi_{k+1}$ is a bijection.
\end{proof}

\subsection{Coordinate Recovery}

\begin{theorem}[Coordinate Recovery]
\label{thm:recovery}
Given a ternary string $\mathbf{t} = (t_1, \ldots, t_k)$, the cell centre coordinates are:
\begin{align}
\Sk(\mathbf{t}) &= \sum_{\substack{j=1 \\ (j-1) \equiv 0 \bmod 3}}^{k} \frac{t_j + 0.5}{3^{\lceil j/3 \rceil}} \label{eq:recovery_sk} \\
\St(\mathbf{t}) &= \sum_{\substack{j=1 \\ (j-1) \equiv 1 \bmod 3}}^{k} \frac{t_j + 0.5}{3^{\lceil j/3 \rceil}} \label{eq:recovery_st} \\
\Se(\mathbf{t}) &= \sum_{\substack{j=1 \\ (j-1) \equiv 2 \bmod 3}}^{k} \frac{t_j + 0.5}{3^{\lceil j/3 \rceil}} \label{eq:recovery_se}
\end{align}
where the $+0.5$ offset selects the midpoint of each ternary interval.
\end{theorem}

\begin{proof}
The $n$-th refinement of dimension $d$ contributes the interval $[t/3^n, (t+1)/3^n)$ for trit value $t$, with midpoint $(t + 0.5)/3^n$. The cell centre is the sum of these midpoints across all refinements for each dimension. The ceiling function $\lceil j/3 \rceil$ counts the refinement depth for the dimension at position $j$.
\end{proof}

\begin{figure*}[!htbp]
    \centering
    \includegraphics[width=\textwidth]{figures/panel_2_ternary_encoding.png}
    \caption{\textbf{Ternary Encoding and Trie Structure.}
    (\textbf{A})~Three-dimensional visualisation of the ternary trie at depth 3. The root node at the top branches into three children per level, producing $3^3 = 27$ leaf cells at depth 3. Only 9 of the 27 cells are occupied (shown as spheres sized by compound count). The largest cells (\texttt{211} with 11 compounds, \texttt{212} with 7) contain the polyatomic-rich clusters, while smaller cells (\texttt{020} with 2, \texttt{110} with 3) capture specific chemical families (heavy halogens, hydrogen halides). Unoccupied cells appear as absent leaves.
    (\textbf{B})~Ternary address heatmap: 39 rows (compounds, sorted lexicographically by trit string) $\times$ 18 columns (trit positions). Cell colour encodes trit value: 0 (blue), 1 (green), 2 (red). Each row is the complete 18-trit ternary address of one compound. Compounds with shared prefixes appear in adjacent rows, revealing the hierarchical clustering structure. Blocks of uniform colour in the leftmost columns correspond to the depth-3 cells from panel~A.
    (\textbf{C})~Resolution curve: number of unique occupied cells versus trit depth (1 through 18). The curve rises from 1 (depth 0, all compounds in one cell) through 9 (depth 3) and 27 (depth 6) to 39 (depth 12, marked by vertical grey line), where all compounds are uniquely resolved. Beyond depth 12, the curve plateaus at 39---additional trits provide finer spatial resolution but no new discriminations for this 39-compound database.
    (\textbf{D})~Cell occupancy at trit depth 3: bar chart showing the number of compounds per occupied cell prefix. Cells are ordered by occupancy. The \texttt{211} cell (11 compounds) dominates, capturing the large cluster of moderate-$\Sk$, moderate-$\St$, moderate-$\Se$ polyatomics. Bar colour indicates the dominant molecular type in each cell: blue$=$diatomic-dominated, green$=$triatomic-dominated, red$=$polyatomic-dominated.}
    \label{fig:ternary_encoding}
\end{figure*}

\subsection{Distance Preservation}

\begin{theorem}[Distance Preservation]
\label{thm:distance}
For two ternary strings $\mathbf{t}$, $\mathbf{t}'$ that first differ at position $j$:
\begin{equation}
d(\varphi(\mathbf{t}), \varphi(\mathbf{t}')) \leq \sqrt{3} \cdot 3^{-\lfloor j/3 \rfloor}.
\label{eq:distance_bound}
\end{equation}
Strings sharing a $k$-trit prefix address cells within distance $\sqrt{3} \cdot 3^{-\lfloor k/3 \rfloor}$.
\end{theorem}

\begin{proof}
If two strings agree on their first $k$ trits, the corresponding points lie in the same level-$k$ cell. After $k$ trits, each dimension has been refined $\lfloor k/3 \rfloor$ or $\lceil k/3 \rceil$ times. The side length along each dimension is at most $3^{-\lfloor k/3 \rfloor}$. The maximum Euclidean distance within the cell is the diagonal:
\begin{equation}
d_{\max} = \sqrt{3 \cdot (3^{-\lfloor k/3 \rfloor})^2} = \sqrt{3} \cdot 3^{-\lfloor k/3 \rfloor}.
\end{equation}
\end{proof}

\begin{corollary}[Prefix Matching is Fuzzy Search]
\label{cor:prefix_fuzzy}
A common prefix of length $k$ guarantees spatial proximity at resolution $3^{-\lfloor k/3 \rfloor}$ per dimension. Truncating the query to $k$ trits performs fuzzy search at this resolution.
\end{corollary}

\subsection{Continuous Emergence}

\begin{theorem}[Continuous Emergence]
\label{thm:continuous}
As $k \to \infty$, the ternary expansion converges to a unique point:
\begin{equation}
\lim_{k \to \infty} C_k(t_1, \ldots, t_k) = \mathbf{S} \in [0,1]^3.
\label{eq:convergence}
\end{equation}
The discrete cell structure converges to the continuous topology. An infinite ternary string specifies a unique point in the continuum.
\end{theorem}

\begin{proof}
The cells $\{C_k\}_{k=0}^\infty$ form a nested decreasing sequence: $C_{k+1} \subset C_k$. Each cell is a closed set. The diameter of $C_k$ decreases as $\sqrt{3} \cdot 3^{-\lfloor k/3 \rfloor} \to 0$ as $k \to \infty$. By the nested closed set theorem in a complete metric space, the intersection $\bigcap_{k=0}^\infty C_k$ contains exactly one point.
\end{proof}

\subsection{Position-Trajectory Duality}

\begin{theorem}[Position-Trajectory Duality]
\label{thm:duality}
Every ternary string simultaneously encodes:
\begin{enumerate}[nosep]
\item A \textbf{position}: the cell (or limiting point) it addresses in $\Sspace$.
\item A \textbf{trajectory}: the path of successive refinements through the partition hierarchy to reach that position.
\end{enumerate}
These are the same mathematical object. The address IS the path.
\end{theorem}

\begin{proof}
The trit string $(t_1, t_2, \ldots, t_k)$ defines the position $\varphi_k(\mathbf{t}) \in \Sspace$ (Theorem~\ref{thm:bijection}) and the trajectory $\Sspace \supset C_1(t_1) \supset C_2(t_1, t_2) \supset \cdots \supset C_k(\mathbf{t})$ (a sequence of nested cells). These are not two different things---the position is the endpoint of the trajectory, and the trajectory is the constructive definition of the position. The data structure (address) and the algorithm (path) are the same ternary string. This dissolves the von Neumann separation between data and instruction at the representational level \citep{vonneumann1945}.
\end{proof}

\begin{remark}
The position-trajectory duality has a practical consequence for database operations. To insert a compound, one follows the trajectory defined by its trit string---the insertion algorithm IS the address. To search for a compound, one follows the same trajectory---the search algorithm IS the address. To compare two compounds, one finds where their trajectories diverge---the comparison IS the common prefix. In each case, the same ternary string serves as data, algorithm, and result. There is no separate ``indexing'' step: the encoding is the index.
\end{remark}



%==============================================================================
% PART II: THE CATEGORICAL COMPOUND DATABASE
%==============================================================================

\part{The Categorical Compound Database}

%==============================================================================
\section{Database Architecture}
\label{sec:architecture}
%==============================================================================

\subsection{The Ternary Trie}

\begin{definition}[Ternary Trie]
\label{def:trie}
A ternary trie $\Trie = (V, E)$ is a rooted tree where:
\begin{itemize}[nosep]
\item The root represents the entire S-entropy space $\Sspace = [0,1]^3$.
\item Each internal node has exactly 3 children, labelled by trit values $\{0, 1, 2\}$.
\item Compounds are stored at leaf nodes (or at the depth to which their trit strings are computed).
\item The depth of a node equals the number of trit positions in its address.
\end{itemize}
\end{definition}

The trie supports three fundamental operations:

\begin{enumerate}
\item \textbf{Insert}$(\mathbf{t}, \text{compound})$: Traverse the trie following the trit string $\mathbf{t} = (t_1, \ldots, t_k)$, creating nodes as needed. Store the compound at the leaf. \textbf{Cost:} $O(k)$.

\item \textbf{SearchExact}$(\mathbf{t})$: Traverse the trie following $\mathbf{t}$. Return compounds stored at the terminal node. \textbf{Cost:} $O(k)$.

\item \textbf{SearchPrefix}$(p)$: Traverse the trie to depth $|p|$ following the prefix $p$. Return ALL compounds in the subtree rooted at that node. \textbf{Cost:} $O(|p| + m)$ where $m$ is the number of matches.
\end{enumerate}

\begin{remark}[Memory Complexity]
The trie has at most $\sum_{j=0}^{k} 3^j = (3^{k+1} - 1)/2$ nodes at depth $k$. However, only the nodes along paths to stored compounds are instantiated. For a database of $N$ compounds at depth $k$, the maximum number of nodes is $N \times k$ (each compound contributes at most $k$ new nodes along its path). In practice, shared prefixes reduce this significantly: at depth 12 with 39 compounds, the trie contains far fewer than $39 \times 12 = 468$ nodes because many compounds share early prefixes.
\end{remark}

\begin{remark}[Concurrent Access]
The trie is a read-mostly data structure: once built, queries require only read access and can be performed concurrently without locking. Insertions add nodes to leaf positions and do not modify existing nodes, enabling lock-free concurrent insertion using compare-and-swap operations on the child pointers. This makes the trie naturally suited to parallel query processing.
\end{remark}

\subsection{Fuzzy Search as Prefix Matching}
\label{sec:fuzzy_search}

\begin{theorem}[Fuzzy Search Completeness]
\label{thm:fuzzy}
For any query point $\mathbf{q} \in \Sspace$ and any radius $\varepsilon > 0$, there exists a trit depth $k$ such that $\mathrm{SearchPrefix}$ at depth $k$ returns all compounds within distance $\varepsilon$ of $\mathbf{q}$, with $k = O(\log_3(1/\varepsilon))$.
\end{theorem}

\begin{proof}
At depth $k$, cells have diameter at most $\sqrt{3} \cdot 3^{-\lfloor k/3 \rfloor}$ (Theorem~\ref{thm:distance}). Set $k = 3 \cdot \lceil \log_3(\sqrt{3}/\varepsilon) \rceil$. Then the cell diameter satisfies:
\begin{equation}
\sqrt{3} \cdot 3^{-\lfloor k/3 \rfloor} \leq \sqrt{3} \cdot 3^{-\log_3(\sqrt{3}/\varepsilon)} = \sqrt{3} \cdot \frac{\varepsilon}{\sqrt{3}} = \varepsilon.
\end{equation}
Any compound within distance $\varepsilon$ of $\mathbf{q}$ lies in the same cell or an adjacent cell at this depth. The query cell and its neighbours can be retrieved by searching a constant number of prefixes.
\end{proof}

\begin{corollary}[Adjustable Resolution]
\label{cor:adjustable}
The resolution of the fuzzy search is continuously adjustable by trit depth:
\begin{itemize}[nosep]
\item Depth 3: coarse resolution, $3^3 = 27$ possible cells, large fuzzy groups.
\item Depth 6: moderate resolution, $3^6 = 729$ possible cells, many compounds unique.
\item Depth 12: fine resolution, $3^{12} = 531{,}441$ possible cells, all 39 compounds individually resolved.
\end{itemize}
\end{corollary}

\begin{figure*}[!htbp]
    \centering
    \includegraphics[width=\textwidth]{figures/panel_3_fuzzy_search.png}
    \caption{\textbf{Fuzzy Search and Chemical Similarity.}
    (\textbf{A})~Fuzzy search demonstration for query $\mathrm{H}_2\mathrm{O}$ at trit depth 3. All 39 compounds are plotted in S-entropy space (small grey dots), with the depth-3 cell containing $\mathrm{H}_2\mathrm{O}$ (prefix \texttt{202}) rendered as a translucent green cube. The six compounds sharing this prefix---$\mathrm{O}_3$, $\mathrm{NO}_2$, $\mathrm{PH}_3$, $\mathrm{CH}_4$, $\mathrm{SO}_2$, $\mathrm{H}_2\mathrm{O}$---appear as large coloured markers inside the cube. All share high $\Sk$, low $\St$, and high $\Se$: spectrally complex molecules with compact frequency ranges and dense harmonic networks.
    (\textbf{B})~Search narrowing cascade for query $\mathrm{CH}_4$, showing matched compounds at four resolutions. At depth 3: 6 matches (green blocks). At depth 6: 4 matches ($\mathrm{NO}_2$, $\mathrm{PH}_3$, $\mathrm{CH}_4$, $\mathrm{SO}_2$; orange blocks). At depth 9: 2 matches ($\mathrm{PH}_3$, $\mathrm{CH}_4$; orange blocks). At depth 12: 1 match ($\mathrm{CH}_4$ uniquely resolved). Each additional 3 trits (one full refinement cycle across all three dimensions) narrows the search, demonstrating adjustable-resolution fuzzy matching without parameter tuning.
    (\textbf{C})~Pairwise ternary similarity matrix for all 39 compounds ($39 \times 39$). Compounds are sorted by ternary address, placing similar compounds adjacent. Colour encodes common prefix length (yellow$=$high similarity, dark purple$=$low similarity). The bright diagonal confirms self-identity. Off-diagonal bright blocks reveal natural clusters: hydrogen halides (HCl--HBr--HI), linear triatomics (CO$_2$--CS$_2$--OCS), and bent triatomics (H$_2$O--SO$_2$--NO$_2$--O$_3$). These clusters emerge from ternary proximity without any chemical knowledge in the encoding.
    (\textbf{D})~Chemical group cohesion analysis. For each of six chemically defined families, paired bars compare mean intra-group ternary similarity (blue) with mean inter-group similarity (grey). Intra$>$inter (blue bar taller) indicates the group is cohesive in ternary space. Five of six families pass: hydrogen halides ($R = 4.38$, strongest), bent triatomics ($R = 3.64$), linear triatomics ($R = 2.44$), halomethanes ($R = 1.77$), homonuclear diatomics ($R = 1.09$). Small hydrocarbons are marginal ($R = 0.98$), reflecting genuine spectral heterogeneity between $\mathrm{CH}_4$, $\mathrm{C}_2\mathrm{H}_2$, $\mathrm{C}_2\mathrm{H}_4$, and $\mathrm{C}_2\mathrm{H}_6$.}
    \label{fig:fuzzy_search}
\end{figure*}

\subsection{Similarity as Common Prefix Length}

\begin{definition}[Ternary Similarity]
\label{def:similarity}
The ternary similarity of compounds $A$, $B$ with trit strings $\mathbf{t}_A$, $\mathbf{t}_B$ is:
\begin{equation}
\mathrm{sim}(A, B) = \max\{k : \mathbf{t}_A[1..k] = \mathbf{t}_B[1..k]\}
\label{eq:similarity}
\end{equation}
the length of their longest common prefix.
\end{definition}

\begin{proposition}[Ultrametric Property]
\label{prop:ultrametric}
The ternary similarity satisfies:
\begin{enumerate}[nosep]
\item $\mathrm{sim}(A, A) = k_{\max}$ \hfill (identity)
\item $\mathrm{sim}(A, B) = \mathrm{sim}(B, A)$ \hfill (symmetry)
\item $\mathrm{sim}(A, B) \geq \min(\mathrm{sim}(A, C), \mathrm{sim}(B, C))$ \hfill (ultrametric inequality)
\end{enumerate}
\end{proposition}

\begin{proof}
Properties (1) and (2) are immediate. For (3): suppose $\mathrm{sim}(A,C) = m$ and $\mathrm{sim}(B,C) = n$ with $m \leq n$. Then $\mathbf{t}_A$ and $\mathbf{t}_C$ agree on their first $m$ trits but differ at position $m+1$, while $\mathbf{t}_B$ and $\mathbf{t}_C$ agree on their first $n \geq m$ trits. In a trie, $A$ and $C$ share a common ancestor at depth $m$, and $B$ and $C$ share a common ancestor at depth $n \geq m$. Since the trie is a tree, $A$ and $B$ must share a common ancestor at depth at least $m = \min(m, n)$, so $\mathrm{sim}(A, B) \geq m = \min(m, n)$.
\end{proof}

The ultrametric inequality is stronger than the triangle inequality and arises because the trie is a tree structure. It implies that ternary similarity defines an equivalence relation at each depth: all compounds sharing a $k$-trit prefix form a cluster at resolution $k$.

\subsection{Structure Prediction}

\begin{theorem}[Property-Based Retrieval]
\label{thm:property_retrieval}
Given constraints on S-entropy coordinates (e.g., $\Sk \in [a,b]$, $\St \in [c,d]$, $\Se \in [e,f]$), the set of matching ternary prefixes can be enumerated in $O(3^k)$ time for depth $k$, and all compounds matching any prefix are retrieved.
\end{theorem}

\begin{proof}
The constraint $\Sk \in [a,b]$ restricts the first coordinate. At depth $k$, the $\Sk$-dimension has been refined $\lfloor k/3 \rfloor$ times, producing $3^{\lfloor k/3 \rfloor}$ intervals of width $3^{-\lfloor k/3 \rfloor}$. The number of intervals intersecting $[a,b]$ is at most $\lceil (b-a) \cdot 3^{\lfloor k/3 \rfloor} \rceil + 1$. The total number of valid prefixes is bounded by the product of valid intervals across all three dimensions, and compounds at each valid prefix are retrieved in $O(k)$ time.
\end{proof}

This theorem inverts the traditional paradigm: instead of ``given molecule, compute properties,'' the question becomes ``given property constraints, which molecules match?'' Figure~\ref{fig:structure_prediction} illustrates this capability.



\begin{figure*}[!htbp]
    \centering
    \includegraphics[width=\textwidth]{figures/panel_4_structure_prediction.png}
    \caption{\textbf{Structure Prediction and Complexity Scaling.}
    (\textbf{A})~Property-based structure prediction in three-dimensional S-entropy space. The translucent orange box delineates a property constraint region ($\Sk \in [0.9, 1.0]$, $\St \in [0.0, 0.3]$), representing molecules with high spectral complexity and narrow frequency range. All 39 compounds are plotted: those inside the constraint region appear as large coloured markers (7 matches: $\mathrm{H}_2\mathrm{O}$, $\mathrm{NO}_2$, $\mathrm{O}_3$, $\mathrm{H}_2\mathrm{S}$, $\mathrm{H}_2\mathrm{CO}$, $\mathrm{PH}_3$, $\mathrm{CH}_4$); those outside appear as small grey dots. This demonstrates the inversion of the search paradigm: from ``given molecule, compute properties'' to ``given property constraints, retrieve molecules.''
    (\textbf{B})~Scaling comparison between brute-force fingerprint search (red line, $O(N \times d)$ with $d = 1024$) and ternary trie search (blue line, $O(k) = O(18)$, constant) on double-logarithmic axes. The $x$-axis spans database sizes from $N = 10$ to $N = 10^6$. The brute-force cost grows linearly with $N$; the trie cost is flat. The vertical dashed line marks the current database ($N = 39$). At PubChem scale ($N = 10^8$, extrapolated), the trie achieves $>10^9\times$ fewer operations per query. The divergence demonstrates that search complexity is fundamentally different when the representation encodes phase space structure directly.
    (\textbf{C})~Ternary distance versus Euclidean distance in S-entropy space for all $\binom{39}{2} = 741$ compound pairs. The $x$-axis shows common ternary prefix length (higher$=$more similar); the $y$-axis shows Euclidean distance $d(\mathbf{S}_1, \mathbf{S}_2)$. The fitted exponential decay curve (red) confirms the Distance Preservation Theorem: longer common prefixes guarantee smaller Euclidean separation. The green horizontal line indicates the noise floor. The correlation validates that ternary prefix matching is a geometrically faithful proxy for proximity in S-entropy space.
    (\textbf{D})~Number of vibrational modes versus evolution entropy $\Se$ (harmonic network density). Point size scales with molecular mass; colour encodes type. Most compounds follow the general trend of increasing $\Se$ with mode count (more modes $\to$ more pairs $\to$ more harmonic coincidences). The outlier $\mathrm{H}_2\mathrm{O}$ (3 modes, $\Se = 1.0$) achieves maximum harmonic density with only three modes because all three pairwise frequency ratios happen to approximate low-order rationals. This demonstrates that $\Se$ captures intrinsic harmonic structure, not mere mode count.}
    \label{fig:structure_prediction}
\end{figure*}
%==============================================================================
\section{Encoding Protocol}
\label{sec:protocol}
%==============================================================================

\subsection{From Vibrational Spectrum to S-Entropy Coordinates}

\begin{protocol}[Molecular Encoding]
\label{prot:encoding}
Given a molecule, its categorical address is obtained as follows:
\begin{enumerate}[nosep]
\item Obtain vibrational frequencies $\{\omega_1, \ldots, \omega_N\}$ from IR/Raman spectroscopy or computational databases (e.g., NIST CCCBDB \citep{johnson2022}).
\item Compute $\Sk$: Shannon entropy of the normalised frequency distribution (Eq.~\ref{eq:s_k_poly} or \ref{eq:s_k_diatomic}).
\item Compute $\St$: logarithmic ratio of the frequency range to the reference range (Eq.~\ref{eq:s_t_poly} or \ref{eq:s_t_diatomic}).
\item Compute $\Se$: harmonic edge density---fraction of mode pairs with rational frequency ratios (Eq.~\ref{eq:s_e}).
\item Generate the interleaved ternary string by Algorithm~1.
\item Insert into the ternary trie.
\end{enumerate}
\end{protocol}

\subsection{Oscillation Counting Interpretation}

Each step of the ternary encoding is an \emph{oscillation-counting observation}: at step $j$, we observe which third of the remaining interval (for dimension $(j-1) \bmod 3$) the coordinate value falls in. This is equivalent to counting three oscillations and recording which phase bin receives the count:

\begin{itemize}[nosep]
\item Step $j$: divide the current interval into three equal parts.
\item Observe which part contains the coordinate value.
\item Record the trit (0, 1, or 2).
\item Refine the interval to the selected part.
\end{itemize}

Thus $k$ trits correspond to $k$ oscillation-counting observations. No algorithmic computation is performed between the spectral measurement and the address generation: no fingerprints are computed, no similarity functions are evaluated, no hash tables are consulted.

The information content per observation is exactly $\log_2 3 \approx 1.585$ bits. After $k$ observations, the molecule has been localised to one of $3^k$ cells, and its address encodes $k \log_2 3$ bits of positional information. For $k = 12$ (the depth at which all 39 compounds are uniquely resolved), this is $12 \times 1.585 \approx 19.0$ bits. For comparison, a 1024-bit Morgan fingerprint encodes 1024 bits but wastes most of them: the effective dimensionality of the fingerprint space for small molecules is far less than 1024 \citep{riniker2013}. The ternary encoding is more efficient because it is matched to the intrinsic dimensionality (3) of the coordinate space.

\begin{theorem}[Computational Equivalence]
\label{thm:computational_equiv}
The ternary encoding of a molecule from its vibrational spectrum requires exactly $k$ oscillation-counting observations for $k$-trit resolution. This is the minimum information-theoretic requirement: $k$ trits encode $\log_2(3^k) = k \cdot \log_2 3 \approx 1.585k$ bits of positional information.
\end{theorem}

\begin{proof}
Each trit selects one of three equal subintervals, conveying $\log_2 3$ bits of information about the coordinate. After $k$ trits, the total information is $k \log_2 3$ bits, localising the point to a cell of volume $3^{-k}$ (or equivalently, distinguishing among $3^k$ possible cells). This is the maximum information extractable from $k$ ternary observations, so the encoding is information-theoretically optimal.
\end{proof}

%==============================================================================
% PART III: VALIDATION
%==============================================================================

\part{Validation}

%==============================================================================
\section{Molecular Test Suite}
\label{sec:test_suite}
%==============================================================================

\subsection{Database Composition}

We validate the categorical compound database on 39 molecules from the NIST Computational Chemistry Comparison and Benchmark Database (CCCBDB) \citep{johnson2022}, spanning:

\begin{itemize}[nosep]
\item 12 diatomics: H$_2$, D$_2$, N$_2$, O$_2$, F$_2$, Cl$_2$, CO, NO, HF, HCl, HBr, HI
\item 10 triatomics: H$_2$O, CO$_2$, SO$_2$, NO$_2$, O$_3$, H$_2$S, HCN, N$_2$O, CS$_2$, OCS
\item 7 tetratomic/tetrahedral: NH$_3$, PH$_3$, CH$_4$, CCl$_4$, SiH$_4$, CF$_4$, H$_2$CO
\item 10 polyatomic: C$_2$H$_2$, C$_2$H$_4$, C$_2$H$_6$, CH$_3$OH, C$_6$H$_6$, CH$_3$F, CH$_3$Cl, CH$_3$Br, HCOOH, CH$_3$CN
\end{itemize}

This selection covers molecules from 2 to 12 atoms, molecular masses from 2 to 154 amu, and vibrational mode counts from 1 to 12, providing a representative cross-section of small-molecule chemical space.

Table~\ref{tab:compounds} presents the complete encoding of all 39 compounds, including vibrational mode counts, S-entropy coordinates, and the first 12 trits of each ternary address.

\begin{longtable}{llccccccl}
\caption{Complete compound encodings for all 39 molecules. $N$ = number of vibrational modes. $\Sk$, $\St$, $\Se$ = S-entropy coordinates. Trits = first 12 digits of the interleaved ternary address.}
\label{tab:compounds} \\
\toprule
Name & Formula & Type & $N$ & Mass & $\Sk$ & $\St$ & $\Se$ & First 12 Trits \\
\midrule
\endfirsthead
\toprule
Name & Formula & Type & $N$ & Mass & $\Sk$ & $\St$ & $\Se$ & First 12 Trits \\
\midrule
\endhead
\midrule
\multicolumn{9}{r}{\textit{Continued on next page}} \\
\bottomrule
\endfoot
\bottomrule
\endlastfoot
% Diatomics
H$_2$    & H$_2$   & diatomic & 1 & 2   & 1.000 & 0.462 & 0.000 & 210210200210 \\
D$_2$    & D$_2$   & diatomic & 1 & 4   & 0.680 & 0.494 & 0.000 & 210010010100 \\
N$_2$    & N$_2$   & diatomic & 1 & 28  & 0.529 & 0.757 & 0.000 & 120100220010 \\
O$_2$    & O$_2$   & diatomic & 1 & 32  & 0.359 & 0.750 & 0.000 & 120000020200 \\
F$_2$    & F$_2$   & diatomic & 1 & 38  & 0.203 & 0.741 & 0.000 & 020100210120 \\
Cl$_2$   & Cl$_2$  & diatomic & 1 & 71  & 0.127 & 0.831 & 0.000 & 020110010110 \\
CO       & CO      & diatomic & 1 & 28  & 0.487 & 0.752 & 0.000 & 120100120000 \\
NO       & NO      & diatomic & 1 & 30  & 0.426 & 0.753 & 0.000 & 120000220100 \\
HF       & HF      & diatomic & 1 & 20  & 0.899 & 0.564 & 0.000 & 210220000000 \\
HCl      & HCl     & diatomic & 1 & 36  & 0.656 & 0.603 & 0.000 & 110220210200 \\
HBr      & HBr     & diatomic & 1 & 81  & 0.582 & 0.612 & 0.000 & 110220010210 \\
HI       & HI      & diatomic & 1 & 128 & 0.507 & 0.627 & 0.000 & 110120110220 \\
\midrule
% Triatomics
H$_2$O   & H$_2$O  & triatomic & 3 & 18  & 0.944 & 0.285 & 1.000 & 202222112122 \\
CO$_2$   & CO$_2$  & triatomic & 3 & 44  & 0.897 & 0.419 & 0.667 & 212200020000 \\
SO$_2$   & SO$_2$  & triatomic & 3 & 64  & 0.937 & 0.322 & 0.667 & 202220120020 \\
NO$_2$   & NO$_2$  & triatomic & 3 & 46  & 0.959 & 0.256 & 0.667 & 202220100220 \\
O$_3$    & O$_3$   & triatomic & 3 & 48  & 0.983 & 0.151 & 0.667 & 202210210100 \\
H$_2$S   & H$_2$S  & triatomic & 3 & 34  & 0.950 & 0.265 & 0.333 & 201220110100 \\
HCN      & HCN     & triatomic & 3 & 27  & 0.864 & 0.512 & 0.333 & 211110210120 \\
N$_2$O   & N$_2$O  & triatomic & 3 & 44  & 0.887 & 0.442 & 0.333 & 211100220220 \\
CS$_2$   & CS$_2$  & triatomic & 3 & 76  & 0.861 & 0.450 & 0.667 & 212110200000 \\
OCS      & OCS     & triatomic & 3 & 60  & 0.855 & 0.458 & 0.667 & 212110200010 \\
\midrule
% Tetratomic/tetrahedral
NH$_3$   & NH$_3$  & tetra & 4 & 17  & 0.918 & 0.429 & 0.500 & 211201021211 \\
PH$_3$   & PH$_3$  & tetra & 4 & 34  & 0.947 & 0.284 & 0.667 & 202220110110 \\
CH$_4$   & CH$_4$  & tetra & 4 & 16  & 0.953 & 0.279 & 0.667 & 202220110210 \\
CCl$_4$  & CCl$_4$ & tetra & 4 & 154 & 0.907 & 0.429 & 0.667 & 212200020110 \\
SiH$_4$  & SiH$_4$ & tetra & 4 & 32  & 0.928 & 0.335 & 0.333 & 211200100000 \\
CF$_4$   & CF$_4$  & tetra & 4 & 88  & 0.945 & 0.360 & 0.667 & 212200100120 \\
H$_2$CO  & H$_2$CO & poly  & 6 & 30  & 0.964 & 0.296 & 0.533 & 201221222001 \\
\midrule
% Polyatomic
C$_2$H$_2$  & C$_2$H$_2$  & poly & 5  & 26  & 0.879 & 0.568 & 0.500 & 211121201211 \\
C$_2$H$_4$  & C$_2$H$_4$  & poly & 12 & 28  & 0.949 & 0.441 & 0.606 & 211202121121 \\
C$_2$H$_6$  & C$_2$H$_6$  & poly & 10 & 30  & 0.954 & 0.429 & 0.667 & 212200120210 \\
CH$_3$OH    & CH$_3$OH    & poly & 10 & 32  & 0.953 & 0.423 & 0.622 & 211202121212 \\
C$_6$H$_6$  & C$_6$H$_6$  & poly & 5  & 78  & 0.913 & 0.504 & 0.600 & 211212011110 \\
CH$_3$F     & CH$_3$F     & poly & 6  & 34  & 0.950 & 0.350 & 0.800 & 212201100111 \\
CH$_3$Cl    & CH$_3$Cl    & poly & 6  & 50  & 0.928 & 0.474 & 0.533 & 211211102021 \\
CH$_3$Br    & CH$_3$Br    & poly & 6  & 95  & 0.917 & 0.536 & 0.600 & 211212021210 \\
HCOOH       & HCOOH       & poly & 8  & 46  & 0.930 & 0.580 & 0.607 & 211222101011 \\
CH$_3$CN    & CH$_3$CN    & poly & 8  & 41  & 0.926 & 0.705 & 0.536 & 221201012201 \\
\end{longtable}

\subsection{S-Entropy Coordinate Distribution}

Figure~\ref{fig:sentropy_space} shows the distribution of all 39 compounds in the three-dimensional S-entropy space.



The S-entropy coordinates reveal systematic patterns across molecular types:

\textbf{Knowledge entropy $\Sk$.} The distribution is bimodal. Diatomics span the range $[0.13, 1.00]$, reflecting the wide variation in fundamental frequencies from Cl$_2$ (560 cm$^{-1}$, $\Sk = 0.127$) to H$_2$ (4401 cm$^{-1}$, $\Sk = 1.000$). Polyatomics cluster near $[0.85, 0.98]$ because Shannon entropy is highest when many modes share the frequency range relatively uniformly.

\textbf{Temporal entropy $\St$.} The distribution is broad. Diatomics cluster at $[0.46, 0.83]$ because the vibrational-to-rotational frequency ratio is typically large ($10^2$--$10^3$). Polyatomics span $[0.15, 0.70]$, with smaller molecules having lower $\St$ (narrower frequency ranges) and larger molecules having higher $\St$.

\textbf{Evolution entropy $\Se$.} This coordinate most clearly separates molecular types. All diatomics have $\Se = 0$ (no mode pairs). Triatomics and polyatomics range from $\Se = 0.33$ (few harmonic coincidences) to $\Se = 1.00$ (H$_2$O, where all pairs are harmonic). The boundary at $\Se = 0$ sharply distinguishes diatomics from polyatomics without any structural information being encoded.

\textbf{Coordinate correlations.} The three coordinates show weak but structured correlations. Molecules with high $\Sk$ (many uniformly distributed modes) tend to have moderate $\Se$ (0.5--0.7), because a uniform frequency distribution implies many near-rational ratios. However, the correlation is imperfect: H$_2$O has both high $\Sk$ (0.944) and maximal $\Se$ (1.000), while SiH$_4$ has high $\Sk$ (0.928) but lower $\Se$ (0.333). The partial independence of the coordinates confirms that they capture genuinely different aspects of molecular oscillatory structure.

%==============================================================================
\section{Fuzzy Search Demonstration}
\label{sec:fuzzy_demo}
%==============================================================================

\subsection{Resolution Cascade}

Figure~\ref{fig:ternary_encoding} illustrates the ternary encoding scheme, showing how successive trits refine the phase space cell.



For five representative query molecules, Table~\ref{tab:fuzzy_cascade} shows how the search results narrow as trit depth increases from 3 to 12.

\begin{table}[H]
\centering
\caption{Fuzzy search results at increasing trit depth. Each entry shows the number of matches and, for small result sets, the matching compounds. At depth 12, all queries resolve to exactly one compound.}
\label{tab:fuzzy_cascade}
\begin{tabular}{lllll}
\toprule
Query & Depth 3 & Depth 6 & Depth 9 & Depth 12 \\
\midrule
H$_2$O & 6: O$_3$, NO$_2$, & 1: H$_2$O & 1: H$_2$O & 1: H$_2$O \\
       & PH$_3$, CH$_4$, SO$_2$, H$_2$O & & & \\[3pt]
CH$_4$ & 6: O$_3$, NO$_2$, & 4: NO$_2$, PH$_3$, & 2: PH$_3$, CH$_4$ & 1: CH$_4$ \\
       & PH$_3$, CH$_4$, SO$_2$, H$_2$O & CH$_4$, SO$_2$ & & \\[3pt]
CO$_2$ & 7: CS$_2$, OCS, CO$_2$, & 4: CO$_2$, CCl$_4$, & 2: CO$_2$, CCl$_4$ & 1: CO$_2$ \\
       & CCl$_4$, CF$_4$, C$_2$H$_6$, CH$_3$F & CF$_4$, C$_2$H$_6$ & & \\[3pt]
C$_2$H$_6$ & 7: (same cell as CO$_2$) & 4: CO$_2$, CCl$_4$, & 1: C$_2$H$_6$ & 1: C$_2$H$_6$ \\
       & & CF$_4$, C$_2$H$_6$ & & \\[3pt]
HCl    & 3: HI, HBr, HCl & 2: HBr, HCl & 1: HCl & 1: HCl \\
\bottomrule
\end{tabular}
\end{table}

The most striking result is HCl at depth 3. The three-trit prefix \texttt{110} returns exactly $\{\text{HI}, \text{HBr}, \text{HCl}\}$---the hydrogen halides. No molecular fingerprints were computed. No Tanimoto coefficients were evaluated. No SMARTS patterns were matched. The ternary phase space address alone groups these chemically related compounds, because they share similar S-entropy coordinates: moderate $\Sk$ (single mode at moderate frequency), moderate $\St$ (similar vibrational-to-rotational ratios), and $\Se = 0$ (diatomics).

The resolution cascade for CH$_4$ is also instructive. At depth 3 (prefix \texttt{202}), CH$_4$ shares a cell with O$_3$, NO$_2$, PH$_3$, SO$_2$, and H$_2$O---six molecules with high $\Sk$, low $\St$, and high $\Se$. At depth 6 (prefix \texttt{202220}), the cell narrows to NO$_2$, PH$_3$, CH$_4$, and SO$_2$---molecules with very high $\Sk$ and moderate $\Se$. At depth 9 (prefix \texttt{202220110}), only PH$_3$ and CH$_4$ remain: both tetrahedral hydrides with nearly identical S-entropy profiles. At depth 12, CH$_4$ is uniquely resolved. This progressive narrowing---from ``high-symmetry molecules'' to ``tetrahedral hydrides'' to ``methane specifically''---is exactly the multi-scale search behaviour the framework is designed to provide.

\subsection{Nearest Neighbours}

Figure~\ref{fig:fuzzy_search} illustrates the fuzzy search mechanism, showing how prefix truncation produces progressively broader search results.



Table~\ref{tab:nearest} shows the nearest neighbours of selected query compounds, ranked by ternary similarity (longest common prefix length).

\begin{table}[H]
\centering
\caption{Nearest neighbours by ternary similarity (common prefix length). Higher similarity indicates more shared trit positions.}
\label{tab:nearest}
\begin{tabular}{lllll}
\toprule
Query & 1st Neighbour & 2nd Neighbour & 3rd Neighbour \\
\midrule
H$_2$O & SO$_2$ (5) & NO$_2$ (5) & PH$_3$ (5) \\
CH$_4$ & PH$_3$ (9) & SO$_2$ (7) & NO$_2$ (7) \\
HCl   & HBr (6) & HI (3) & N$_2$ (1) \\
CO$_2$ & CCl$_4$ (9) & CF$_4$ (6) & C$_2$H$_6$ (6) \\
CH$_3$Cl & C$_6$H$_6$ (5) & CH$_3$Br (5) & NH$_3$ (4) \\
C$_2$H$_6$ & CF$_4$ (7) & CO$_2$ (6) & CCl$_4$ (6) \\
\bottomrule
\end{tabular}
\end{table}

CH$_4$ and PH$_3$ share a 9-trit prefix---the longest common prefix among non-identical compounds in the database. Both are tetrahedral hydrides with similar vibrational structure: four modes spanning comparable frequency ranges, with identical harmonic connectivity ($\Se = 0.667$). The framework discovers this structural relationship from vibrational frequencies alone, without being programmed to recognise molecular symmetry or bonding patterns.

The CO$_2$--CCl$_4$ pairing (9 trits shared) is initially surprising: these molecules have very different chemistry. However, they share remarkably similar S-entropy profiles: both have $\Sk \approx 0.90$ (relatively uniform frequency distribution), $\St \approx 0.42$ (similar frequency range relative to reference), and $\Se = 0.667$ (same harmonic connectivity fraction). The encoding captures spectral similarity, not necessarily chemical similarity, and CO$_2$ and CCl$_4$ happen to have spectrally analogous vibrational structures despite their chemical differences.

%==============================================================================
\section{Chemical Similarity Validation}
\label{sec:similarity}
%==============================================================================

\subsection{Validation Methodology}

The central question of this section is: does the ternary encoding, derived purely from vibrational frequencies, reproduce known chemical relationships? If it does, this validates both the S-entropy coordinate system and the ternary encoding scheme. If it does not, this identifies limitations of the spectral representation.

We emphasise that no chemical knowledge is encoded in the S-entropy coordinates. The encoding uses only vibrational frequencies, a harmonic proximity threshold, and reference frequency bounds. It does not know about bond types, functional groups, electronegativity, or molecular symmetry. Any chemical grouping that emerges from the ternary encoding is a discovery, not a tautology.

\subsection{Chemical Families}

We define six chemically meaningful families to test whether ternary proximity reproduces known chemical groupings:

\begin{enumerate}[nosep]
\item \textbf{Halomethanes:} CH$_3$F, CH$_3$Cl, CH$_3$Br --- substituted methanes with halogen substituents.
\item \textbf{Hydrogen halides:} HF, HCl, HBr, HI --- diatomic hydrides of halogens.
\item \textbf{Homonuclear diatomics:} H$_2$, D$_2$, N$_2$, O$_2$, F$_2$, Cl$_2$ --- elemental diatomic molecules.
\item \textbf{Small hydrocarbons:} CH$_4$, C$_2$H$_2$, C$_2$H$_4$, C$_2$H$_6$ --- saturated and unsaturated hydrocarbons.
\item \textbf{Linear triatomics:} CO$_2$, CS$_2$, N$_2$O, OCS --- linear molecules with three atoms.
\item \textbf{Bent triatomics:} H$_2$O, SO$_2$, NO$_2$, O$_3$, H$_2$S --- nonlinear molecules with three atoms.
\end{enumerate}

\subsection{Cohesion Metric}

\begin{definition}[Cohesion Ratio]
\label{def:cohesion}
For a chemical group $G$ with members $\{g_1, \ldots, g_m\}$ in a database of all compounds $\{c_1, \ldots, c_N\}$:
\begin{align}
\text{Intra-group similarity} &= \frac{1}{\binom{m}{2}} \sum_{i < j} \mathrm{sim}(g_i, g_j) \label{eq:intra} \\
\text{Inter-group similarity} &= \frac{1}{m(N-m)} \sum_{g \in G} \sum_{c \notin G} \mathrm{sim}(g, c) \label{eq:inter} \\
R &= \frac{\text{Intra-group similarity}}{\text{Inter-group similarity}} \label{eq:cohesion_ratio}
\end{align}
$R > 1$ indicates the group is more internally cohesive than externally---members are closer to each other in ternary space than to non-members.
\end{definition}

\subsection{Results}

Table~\ref{tab:validation} presents the cohesion analysis for all six chemical families.

\begin{table}[H]
\centering
\caption{Chemical group validation results. Intra = mean pairwise ternary similarity within the group. Inter = mean ternary similarity between group members and non-members. $R$ = intra/inter ratio. $R > 1$ indicates ternary cohesion.}
\label{tab:validation}
\begin{tabular}{lccccc}
\toprule
Family & Members & Intra & Inter & $R$ & Status \\
\midrule
Halomethanes         & 3 & 3.00 & 1.69 & 1.77 & PASS \\
Hydrogen halides     & 4 & 2.00 & 0.46 & 4.38 & PASS \\
Homonuclear diatomics & 6 & 0.67 & 0.61 & 1.09 & PASS \\
Small hydrocarbons   & 4 & 1.67 & 1.71 & 0.98 & MARGINAL \\
Linear triatomics    & 4 & 3.67 & 1.50 & 2.44 & PASS \\
Bent triatomics      & 5 & 3.70 & 1.02 & 3.64 & PASS \\
\bottomrule
\end{tabular}
\end{table}

\subsection{Discussion of Results}

Five of six families show $R > 1.0$, indicating that ternary proximity successfully reproduces chemical groupings without any chemical knowledge being encoded in the S-entropy coordinates.

\textbf{Hydrogen halides ($R = 4.38$):} The strongest cohesion. All four members have $\Se = 0$ (single mode, no pairs), similar $\St$ values (moderate vibrational-to-rotational ratios), and $\Sk$ values that decrease systematically with increasing halogen mass (HF: 0.899, HCl: 0.656, HBr: 0.582, HI: 0.507). The common three-trit prefix \texttt{110} captures three of the four members (HCl, HBr, HI); HF has a higher $\Sk$ that places it in cell \texttt{210} alongside D$_2$ and H$_2$. Within the ternary encoding, the hydrogen halides form a tight spectral cluster that is well-separated from all other compound types.

\textbf{Bent triatomics ($R = 3.64$):} The second strongest cohesion. H$_2$O, SO$_2$, NO$_2$, O$_3$ all share the \texttt{202} prefix at depth 3, reflecting their high $\Sk$ (uniform frequency distribution among three modes), low-to-moderate $\St$ (moderate frequency ranges), and high $\Se$ ($= 0.667$ for most). H$_2$S is the outlier with lower $\Se = 0.333$, but it still clusters with the group at depth 2.

\textbf{Linear triatomics ($R = 2.44$):} CO$_2$, CS$_2$, OCS share the \texttt{212} prefix, and CS$_2$--OCS share a 7-trit prefix. N$_2$O has lower $\Se$ (0.333 vs 0.667) but still clusters with the group at depth 3 (\texttt{211}).

\textbf{Halomethanes ($R = 1.77$):} CH$_3$F, CH$_3$Cl, CH$_3$Br all share the \texttt{211} prefix and cluster within 3--5 trit similarity. Their S-entropy profiles reflect the common CH$_3$ vibrational framework with systematic variation in the C--X stretching frequency as the halogen mass increases.

\textbf{Homonuclear diatomics ($R = 1.09$):} The weakest passing group. The six members span a wide range of $\Sk$ (from Cl$_2$ at 0.127 to H$_2$ at 1.000) because their fundamental frequencies vary by nearly an order of magnitude. They share $\Se = 0$ but have diverse $\Sk$ and $\St$ values, producing a diffuse cluster that barely exceeds the inter-group baseline.

\textbf{Small hydrocarbons ($R = 0.98$, MARGINAL):} The one group that does not achieve clear cohesion. This is not a failure of the encoding but a reflection of genuine chemical heterogeneity. Methane (CH$_4$, tetrahedral, 4 modes) has a fundamentally different vibrational topology from acetylene (C$_2$H$_2$, linear, 5 modes), ethylene (C$_2$H$_4$, planar, 12 modes), and ethane (C$_2$H$_6$, 10 modes). The label ``small hydrocarbons'' groups molecules with very different oscillatory structures, and the S-entropy encoding correctly reflects this diversity. The encoding reports spectral similarity, not chemical similarity \emph{per se}, and molecules that are chemically related but spectrally dissimilar will not cluster. This is a feature, not a bug: it means the encoding captures a different and potentially more fundamental aspect of molecular identity than traditional chemical classification.

\subsection{Summary of Validation}

The validation demonstrates that the S-entropy encoding captures genuine chemical structure, with the following caveats:

\begin{enumerate}[nosep]
\item Groups defined by vibrational similarity (hydrogen halides, bent triatomics, linear triatomics) show strong ternary cohesion ($R > 2$).
\item Groups defined by structural similarity that correlates with vibrational similarity (halomethanes) show moderate cohesion ($R \approx 1.8$).
\item Groups defined by structural similarity that does NOT correlate with vibrational similarity (small hydrocarbons: CH$_4$ vs C$_2$H$_2$) show marginal or no cohesion ($R \approx 1$).
\item Groups defined by elemental composition alone (homonuclear diatomics) show weak cohesion ($R \approx 1.1$).
\end{enumerate}

This pattern is exactly what one would expect from a spectral encoding. The S-entropy coordinates are blind to chemical identity \emph{per se}; they measure oscillatory structure. When chemical identity correlates with oscillatory structure (as it often does), the encoding reproduces the chemical grouping. When it does not (as for ``small hydrocarbons,'' which span tetrahedral, linear, and planar geometries), the encoding correctly reports the spectral diversity.

\subsection{Emergent Clustering at Depth 3}

At three-trit resolution, the 39 compounds partition into 9 occupied cells (Table~\ref{tab:depth3}). These groupings emerge from the S-entropy encoding alone, without any chemical knowledge.

\begin{table}[H]
\centering
\caption{Clustering at trit depth 3. The 39 compounds occupy 9 cells of the $3^3 = 27$ possible cells. Chemical interpretations are provided post hoc---they were not used in the encoding.}
\label{tab:depth3}
\begin{tabular}{lcp{8.5cm}}
\toprule
Cell & $N$ & Members and Chemical Interpretation \\
\midrule
\texttt{[020]} & 2 & F$_2$, Cl$_2$ --- heavy homonuclear halogens \\
\texttt{[110]} & 3 & HI, HBr, HCl --- hydrogen halides \\
\texttt{[120]} & 4 & O$_2$, NO, CO, N$_2$ --- diatomics with similar $\omega$ and $\Brot$ \\
\texttt{[201]} & 2 & H$_2$S, H$_2$CO --- moderate $\Sk$, low $\St$, low $\Se$ \\
\texttt{[202]} & 6 & O$_3$, NO$_2$, PH$_3$, CH$_4$, SO$_2$, H$_2$O --- high $\Sk$, low $\St$, high $\Se$ \\
\texttt{[210]} & 3 & D$_2$, H$_2$, HF --- light diatomics with high $\omega$ \\
\texttt{[211]} & 11 & N$_2$O, HCN, C$_2$H$_2$, SiH$_4$, NH$_3$, C$_2$H$_4$, CH$_3$OH, CH$_3$Cl, C$_6$H$_6$, CH$_3$Br, HCOOH --- large polyatomic cluster \\
\texttt{[212]} & 7 & CS$_2$, OCS, CO$_2$, CCl$_4$, CF$_4$, C$_2$H$_6$, CH$_3$F --- moderate $\Sk$, moderate $\St$, high $\Se$ \\
\texttt{[221]} & 1 & CH$_3$CN --- high $\Sk$, high $\St$, moderate $\Se$ \\
\bottomrule
\end{tabular}
\end{table}

Several chemically sensible groupings emerge:

\begin{itemize}[nosep]
\item \texttt{[020]}: F$_2$ and Cl$_2$ are both homonuclear halogens with low $\Sk$ (low fundamental frequency) and high $\St$ (large vibrational-to-rotational ratio).
\item \texttt{[110]}: HI, HBr, HCl---exactly the three heavier hydrogen halides---grouped by moderate $\Sk$, moderate $\St$, zero $\Se$.
\item \texttt{[120]}: O$_2$, NO, CO, N$_2$---four diatomics with similar vibrational frequencies (1580--2330 cm$^{-1}$) and rotational constants (1.4--2.0 cm$^{-1}$).
\item \texttt{[210]}: D$_2$, H$_2$, HF---the three lightest diatomics with the highest vibrational frequencies.
\end{itemize}

These groupings arise purely from the oscillatory structure of each molecule, encoded as three numbers and interleaved into a ternary string. No molecular graphs, no bond types, no functional groups, no chemical rules were used.

\subsection{Refinement at Depth 6 and 9}

As the trit depth increases, the clusters refine. At depth 6, 27 cells are occupied, with 8 containing multiple compounds:

\begin{table}[H]
\centering
\caption{Multi-occupancy cells at trit depth 6. Single-occupancy cells (19 of 27) are omitted.}
\label{tab:depth6}
\begin{tabular}{lcp{8cm}}
\toprule
Cell & $N$ & Members \\
\midrule
\texttt{[202220]} & 4 & NO$_2$, PH$_3$, CH$_4$, SO$_2$ \\
\texttt{[212200]} & 4 & CO$_2$, CCl$_4$, CF$_4$, C$_2$H$_6$ \\
\texttt{[110220]} & 2 & HBr, HCl \\
\texttt{[120000]} & 2 & O$_2$, NO \\
\texttt{[120100]} & 2 & CO, N$_2$ \\
\texttt{[211202]} & 2 & C$_2$H$_4$, CH$_3$OH \\
\texttt{[211212]} & 2 & C$_6$H$_6$, CH$_3$Br \\
\texttt{[212110]} & 2 & CS$_2$, OCS \\
\bottomrule
\end{tabular}
\end{table}

The CO--N$_2$ pairing at depth 6 is notable: these molecules are isoelectronic (both 14 electrons, triple bond, similar bond lengths), and their vibrational frequencies are similar (CO: 2143 cm$^{-1}$, N$_2$: 2330 cm$^{-1}$) with comparable rotational constants (CO: 1.93 cm$^{-1}$, N$_2$: 1.99 cm$^{-1}$). The encoding discovers this isoelectronic relationship from spectral data alone.

At depth 9, 35 cells are occupied with only 4 multi-occupancy cells remaining: PH$_3$--CH$_4$ (tetrahedral hydrides), C$_2$H$_4$--CH$_3$OH (coincidentally similar S-entropy profiles), CS$_2$--OCS (linear triatomics with S-containing bonds), and CO$_2$--CCl$_4$ (spectrally analogous despite chemical dissimilarity). By depth 12, all 39 compounds are uniquely resolved: every molecule has a distinct 12-trit address.

%==============================================================================
\section{Structure Prediction Examples}
\label{sec:prediction}
%==============================================================================

The categorical compound database supports \emph{inverse} queries: given constraints on S-entropy coordinates, which molecules match? This inverts the traditional cheminformatics pipeline, which computes properties from structures.

\subsection{Property-Constrained Search}

\begin{example}[High spectral complexity, low temporal spread]
The constraints $\Sk \in [0.9, 1.0]$ and $\St \in [0.0, 0.3]$ select molecules with nearly uniform frequency distributions and narrow frequency ranges. The database returns 7 compounds: H$_2$O, NO$_2$, O$_3$, H$_2$S, H$_2$CO, PH$_3$, CH$_4$. These are all small molecules (3--5 atoms) with compact vibrational spectra---exactly the class one would expect from the constraint.
\end{example}

\begin{example}[Low harmonic connectivity]
The constraint $\Se \in [0.0, 0.15]$ returns exactly the 12 diatomics: H$_2$, D$_2$, N$_2$, O$_2$, F$_2$, Cl$_2$, CO, NO, HF, HCl, HBr, HI. All have $\Se = 0$ because they possess only one vibrational mode and hence no mode pairs. The $\Se$ coordinate alone perfectly separates diatomics from all other molecule types in the database.
\end{example}

\begin{example}[Simple spectrum, wide temporal spread]
The constraints $\Sk \in [0.0, 0.4]$ and $\St \in [0.5, 1.0]$ select molecules with one dominant low-frequency mode and a large vibrational-to-rotational frequency ratio. The database returns 3 compounds: O$_2$, F$_2$, Cl$_2$---heavy homonuclear diatomics with low fundamental frequencies and small rotational constants. These are precisely the molecules whose oscillatory structure is characterised by a single slow vibration spanning many rotational timescales.
\end{example}

These examples demonstrate that the S-entropy coordinates provide chemically interpretable axes. Constraints on $\Sk$ select by spectral complexity, constraints on $\St$ select by timescale hierarchy, and constraints on $\Se$ select by harmonic connectivity. The coordinate system is not merely a mathematical convenience; it captures physically meaningful aspects of molecular oscillatory structure.

%==============================================================================
\section{Complexity Analysis}
\label{sec:complexity}
%==============================================================================

\subsection{Traditional Molecular Search}

In traditional cheminformatics, the standard pipeline for molecular similarity search has the following computational profile:

\begin{itemize}[nosep]
\item \textbf{Fingerprint:} Morgan/ECFP fingerprint \citep{rogers2010}, typically $d = 1024$ bits.
\item \textbf{Similarity:} Tanimoto coefficient $T(A,B) = |A \cap B|/|A \cup B|$, requiring $O(d)$ bit operations.
\item \textbf{Database scan:} For $N$ compounds, $O(N)$ pairwise comparisons.
\item \textbf{Total cost per query:} $O(N \times d)$.
\end{itemize}

For the PubChem database \citep{kim2023}: $N \approx 10^8$, $d = 1024$, giving $\sim 10^{11}$ operations per query.

\subsection{Ternary Trie Search}

In the categorical framework:

\begin{itemize}[nosep]
\item \textbf{Encoding:} $O(k)$ oscillation-counting observations for $k$-trit depth.
\item \textbf{Search:} $O(k)$ trie traversals (follow the trit string through the trie).
\item \textbf{Database size:} Does not appear in the complexity---$N$ is absent.
\item \textbf{Resolution:} $k = 18$ trits gives resolution $3^{-6} \approx 0.0014$ per dimension.
\end{itemize}

\begin{theorem}[Scaling Independence]
\label{thm:scaling}
Ternary trie search requires $O(k)$ operations for $k$-trit resolution, regardless of the number of compounds $N$ in the database.
\end{theorem}

\begin{proof}
The trie traversal follows $k$ edges from root to leaf, one per trit. Each edge traversal is $O(1)$ (array lookup on a 3-element array). The total cost is $O(k)$, independent of $N$. The number of compounds affects only the \emph{storage} (trie size), not the \emph{search time}.

More precisely, let $T(k)$ denote the time to search at depth $k$. Each step consists of:
\begin{enumerate}[nosep]
\item Read the current trit value from the query string: $O(1)$.
\item Follow the corresponding child pointer: $O(1)$.
\item Check whether the child exists (non-null test): $O(1)$.
\end{enumerate}
Thus $T(k) = 3k = O(k)$. The constant factor 3 is independent of $N$, $d$, or any other database parameter.
\end{proof}

\begin{remark}
The $O(k)$ complexity holds for exact search. For prefix search (fuzzy search), the additional cost is $O(m)$ where $m$ is the number of compounds in the subtree. At coarse resolutions (small $k$), $m$ may be large, approaching $N$ in the degenerate case $k = 0$. The choice of $k$ therefore trades off between search precision and result set size, but the \emph{search traversal} itself remains $O(k)$.
\end{remark}

\subsection{Speedup Analysis}

Table~\ref{tab:speedup} compares the operation counts for brute-force fingerprint search against ternary trie lookup at different depths and database sizes.

\begin{table}[H]
\centering
\caption{Speedup of ternary trie search over brute-force fingerprint comparison ($d = 1024$). For each database size $N$ and trit depth $k$, the speedup is $(N \times d) / k$.}
\label{tab:speedup}
\begin{tabular}{lrrrrr}
\toprule
& \multicolumn{5}{c}{Trit Depth $k$} \\
\cmidrule(lr){2-6}
Database Size $N$ & $k=3$ & $k=6$ & $k=9$ & $k=12$ & $k=18$ \\
\midrule
$39$ (this work)        & $13{,}312$ & $6{,}656$ & $4{,}437$ & $3{,}328$ & $2{,}219$ \\
$10^4$                  & $3.4 \times 10^6$ & $1.7 \times 10^6$ & $1.1 \times 10^6$ & $8.5 \times 10^5$ & $5.7 \times 10^5$ \\
$10^6$                  & $3.4 \times 10^8$ & $1.7 \times 10^8$ & $1.1 \times 10^8$ & $8.5 \times 10^7$ & $5.7 \times 10^7$ \\
$10^8$ (PubChem)        & $3.4 \times 10^{10}$ & $1.7 \times 10^{10}$ & $1.1 \times 10^{10}$ & $8.5 \times 10^9$ & $5.7 \times 10^9$ \\
\bottomrule
\end{tabular}
\end{table}

For the current database ($N = 39$, $d = 1024$, $k = 12$):
\begin{equation}
\text{Speedup} = \frac{39 \times 1024}{12} = \frac{39{,}936}{12} = 3{,}328\times.
\end{equation}

At PubChem scale ($N = 10^8$, $d = 1024$, $k = 18$):
\begin{equation}
\text{Speedup} = \frac{10^8 \times 1024}{18} \approx 5.7 \times 10^9 \quad (\text{nearly } 10^{10}\text{-fold}).
\end{equation}

This speedup is not achieved through algorithmic cleverness (parallelism, caching, pre-filtering) but through a change of representation. The ternary trie makes the database size irrelevant to the search cost because the search follows the query's address, not the database's contents. Whether the trie contains 39 nodes or $10^8$ nodes, following a path of length $k$ takes $k$ steps.

It is important to note that the comparison is between the \emph{search} operations, not the total pipeline. Computing the S-entropy coordinates from a vibrational spectrum requires $O(N_{\mathrm{modes}}^2)$ work (dominated by the pairwise harmonic proximity check for $\Se$). For a molecule with 30 modes, this is $\sim 450$ pair evaluations. This one-time encoding cost is comparable to computing a Morgan fingerprint and is amortised across all subsequent queries.

\subsection{Resolution at Each Depth}

Table~\ref{tab:resolution} shows the resolution achieved at each trit depth, the number of possible cells, and the resolution per coordinate dimension.

\begin{table}[H]
\centering
\caption{Resolution and cell structure at each trit depth. Trits per dimension = $\lfloor k/3 \rfloor$. Resolution per axis = $3^{-\lfloor k/3 \rfloor}$.}
\label{tab:resolution}
\begin{tabular}{cccc}
\toprule
Depth $k$ & Trits/Dimension & Total Cells ($3^k$) & Resolution/Axis \\
\midrule
3  & 1 & 27 & 0.3333 \\
6  & 2 & 729 & 0.1111 \\
9  & 3 & 19{,}683 & 0.0370 \\
12 & 4 & 531{,}441 & 0.0123 \\
15 & 5 & 14{,}348{,}907 & 0.0041 \\
18 & 6 & 387{,}420{,}489 & 0.0014 \\
\bottomrule
\end{tabular}
\end{table}

\subsection{Oscillation Counting vs Algorithmic Computation}

The fundamental difference between traditional search and categorical search is not merely quantitative (fewer operations) but qualitative (different computational paradigm).

In traditional search, the computer \emph{computes} a similarity function: it transforms two molecular representations through a series of algorithmic steps (fingerprint generation, bit operations, division) to produce a numerical score. The computation is indirect---the similarity is a derived quantity, not a direct property of the representation.

In categorical search, the computer \emph{counts oscillations} and \emph{reads the address}. No similarity function is evaluated. No fingerprint is compared. The trie encodes the bounded phase space directly, and search is phase space navigation. Each trit is one observation, and the observations collectively specify a location.

\begin{table}[H]
\centering
\caption{Paradigm comparison between traditional and categorical molecular search.}
\label{tab:paradigm}
\begin{tabular}{lll}
\toprule
Aspect & Traditional & Categorical \\
\midrule
Representation & SMILES $\to$ fingerprint & Spectrum $\to$ $(\Sk, \St, \Se)$ $\to$ trits \\
Search & Pairwise comparison $O(Nd)$ & Trie traversal $O(k)$ \\
Similarity & Tanimoto coefficient & Common prefix length \\
Resolution & Fixed (fingerprint length) & Adjustable (trit depth) \\
Fuzzy matching & Threshold on Tanimoto & Prefix truncation \\
Prediction & Screen library or ML & Constraint $\to$ prefix $\to$ trie \\
\bottomrule
\end{tabular}
\end{table}

%==============================================================================
% PART IV: DISCUSSION AND CONCLUSION
%==============================================================================

\part{Discussion and Conclusion}

%==============================================================================
\section{Discussion}
\label{sec:discussion}
%==============================================================================

\subsection{What Has Been Demonstrated}

This paper has established the following:

\begin{enumerate}
\item \textbf{Intrinsic molecular coordinates.} The S-entropy coordinates $(\Sk, \St, \Se)$ provide a three-dimensional coordinate system for molecular identity, derived entirely from vibrational spectra. No molecular graphs, connectivity tables, or structural descriptors are required.

\item \textbf{Ternary addressing.} The interleaved ternary encoding maps molecules to addresses in a trie, where search is $O(k)$ independent of database size. The encoding is natural (base-3 for 3D space), distance-preserving, and hierarchical.

\item \textbf{Emergent chemical similarity.} Chemical families cluster in ternary space without chemical knowledge being encoded. The hydrogen halides, bent triatomics, and linear triatomics show particularly strong ternary cohesion, with groupings emerging at the coarsest resolution (depth 3).

\item \textbf{Structure prediction.} The framework inverts the traditional query: instead of ``given a molecule, find similar ones,'' the question becomes ``given spectral property constraints, which addresses match?'' This is answered by enumerating valid ternary prefixes.
\end{enumerate}

\subsection{The Role of Fuzzy Representation}

A compound at $k$ trits is a \emph{region}, not a point. At low resolution, many compounds share the same cell. This is not a deficiency---it is the mechanism of fuzzy matching.

Traditional cheminformatics achieves fuzzy matching by choosing a Tanimoto threshold: compounds with $T > 0.7$ are ``similar.'' This threshold is arbitrary and must be tuned for each application. In the categorical framework, fuzzy matching is achieved by choosing a trit depth: compounds sharing a $k$-trit prefix are ``similar at resolution $k$.'' The resolution is principled---each additional trit corresponds to one more oscillation-counting observation---and there is no arbitrary threshold to tune.

The hierarchy of resolutions forms a multi-scale view of chemical space. At depth 3, chemistry is seen through a coarse lens: broad families (diatomics, triatomics, polyatomics) are distinguishable. At depth 6, finer distinctions emerge (hydrogen halides vs other diatomics, bent vs linear triatomics). At depth 12, every compound in our database is individually resolved. At depth 18, the resolution is $\sim 0.001$ per axis---sufficient to distinguish molecules whose S-entropy coordinates differ by less than 0.1\%.

\subsection{From Counting to Chemistry}

The encoding protocol (Section~\ref{sec:protocol}) reveals a direct correspondence between computational operations and physical observations:

\begin{itemize}[nosep]
\item Each trit = one refinement = one oscillation-counting observation.
\item The CPU clock IS an oscillator; counting its cycles IS spectroscopic measurement.
\item Oscillation counting = partitioning phase space = categorical measurement.
\end{itemize}

This correspondence is not metaphorical. The ternary expansion of a coordinate value is mathematically identical to the process of repeatedly observing which third of an interval a physical oscillator occupies. The ``computation'' of the ternary address from the spectrum is a series of such observations, each extracting $\log_2 3$ bits of positional information. No algorithmic transformation intervenes between measurement and address.

The triple equivalence---oscillation, counting, partition---means that:

\begin{enumerate}[nosep]
\item Counting oscillations $\equiv$ partitioning phase space $\equiv$ categorical measurement.
\item No algorithmic computation is required between measurement and search.
\item The measurement IS the search.
\end{enumerate}

To make this concrete: when the encoding algorithm computes $t_j = \lfloor 3 \cdot r_{\mathrm{dim}} \rfloor$, it is asking ``in which third of the remaining interval does this coordinate fall?'' This is operationally equivalent to observing a ternary oscillator (one with three distinguishable phases) and recording which phase is occupied. The ``oscillation'' is the cyclic partition of the unit interval into thirds; the ``counting'' is the recording of the trit value; the ``partition'' is the refinement of the phase space cell. These three descriptions are the same operation viewed from different perspectives.

\subsection{The Trie as a Physical Object}

The ternary trie is not merely a data structure in the computer science sense. It is a discretised representation of the bounded phase space $\Sspace = [0,1]^3$. Each node in the trie corresponds to a physical region of S-entropy space---a set of molecules whose oscillatory structures are indistinguishable at that resolution. The edges of the trie correspond to refinement operations---additional oscillation-counting observations that narrow the region.

This physical interpretation has consequences. The trie is not designed; it is \emph{derived} from the coordinate space. Adding a new compound to the database does not change the trie structure (only populates an existing node). Removing a compound does not invalidate any other compound's address. The addresses are stable because they are determined by the physics (vibrational spectrum), not by the database contents.

In traditional databases, adding or removing entries can change the results of similarity searches (e.g., by altering clustering thresholds or nearest-neighbour rankings). In the categorical database, the address of a compound is independent of all other compounds. This is a direct consequence of the intrinsic coordinate system: the molecule's position in S-entropy space is determined by its own vibrational spectrum, not by its relationship to other molecules.

\subsection{Scaling to Large Databases}

The theoretical scaling advantage of $O(k)$ vs $O(Nd)$ becomes practically significant at large $N$. Consider the following concrete scenario. A pharmaceutical company screens a compound library of $N = 10^6$ molecules against a query hit. In the traditional pipeline:

\begin{itemize}[nosep]
\item Compute Morgan fingerprint for the query: $O(1)$ (single molecule).
\item Compare against all $10^6$ database fingerprints: $10^6 \times 1024 \approx 10^9$ bit operations.
\item Sort by Tanimoto score: $O(N \log N) \approx 2 \times 10^7$ comparisons.
\item Total: $\sim 10^9$ operations per query.
\end{itemize}

In the categorical pipeline:

\begin{itemize}[nosep]
\item Compute S-entropy coordinates from measured vibrational spectrum: $O(N_{\mathrm{modes}}^2)$, typically $< 1000$ operations.
\item Generate trit string: $O(k)$, typically 12--18 operations.
\item Search trie: $O(k)$, typically 12--18 operations.
\item Total: $\sim 1000$ operations per query.
\end{itemize}

The categorical pipeline is approximately $10^6$ times faster for this scenario. Moreover, the categorical pipeline produces not just a ranked list but a \emph{hierarchical} result: the set of matches at each resolution level from coarse to fine. This multi-scale output is intrinsic to the trie structure and requires no additional computation.

\subsection{Comparison with Existing Methods}

\textbf{Morgan/ECFP fingerprints} \citep{rogers2010}: These encode molecular substructure as fixed-length binary vectors. They are highly effective for substructure similarity but require $O(N)$ pairwise comparison. The fingerprint is an extrinsic descriptor: two molecules with identical fingerprints may have different structures, and the fingerprint length ($d = 1024$ or 2048) is a design choice unrelated to any physical principle. In contrast, S-entropy coordinates are intrinsic (derived from the physical spectrum) and the trit depth is a resolution parameter with physical meaning.

\textbf{Molecular graphs and canonical indexing} \citep{morgan1965}: Graph-based methods encode molecules as labelled graphs and compare them via canonical indexing algorithms (Morgan algorithm) or graph kernels. These are structurally informative but still require $O(N)$ database scans for similarity search. The categorical framework bypasses the graph representation entirely, working directly with the vibrational spectrum.

\textbf{Locality-sensitive hashing (LSH)} \citep{indyk1998}: LSH achieves sublinear search by hashing similar items to the same bucket with high probability. However, LSH is approximate (it may miss near neighbours), the hash functions must be designed for the specific metric space, and the theoretical guarantees are probabilistic. Ternary trie search is exact within resolution (all compounds in the addressed cell are returned), the encoding is derived from physics (not designed), and the guarantees are deterministic.

\textbf{Space-partitioning trees (kd-trees, ball trees) \citep{knuth1998}:} These classical data structures achieve $O(\log N)$ search in low dimensions but degrade to $O(N)$ in high dimensions. The S-entropy space is three-dimensional---low enough for space-partitioning methods to work well. However, the ternary trie is not merely a kd-tree in disguise: the ternary encoding arises from the 3D structure of the coordinate space and interleaves refinements across dimensions in a way that preserves the physical meaning of each trit. In a kd-tree, the splitting dimension and threshold are chosen adaptively based on the data distribution; in the ternary trie, the splitting is fixed and determined by the coordinate system itself. This means the trie structure is independent of the database contents---a property that kd-trees do not share.

\textbf{Trie data structures \citep{fredkin1960}:} The ternary trie is an instance of the classical trie data structure, specialised to a 3-symbol alphabet. The novelty is not the data structure but the encoding: the interleaved ternary address derived from S-entropy coordinates provides a physically meaningful key that makes prefix matching equivalent to spatial proximity search. Standard tries on arbitrary keys (e.g., text strings) do not have this geometric interpretation.

\subsection{Limitations}

\begin{enumerate}
\item \textbf{Fundamental frequencies only.} The current encoding uses only fundamental vibrational frequencies. Overtones, combination bands, and anharmonic effects are not included. Including these would enrich the spectral information available for $\Se$ computation but would not change the encoding framework.

\item \textbf{Database size.} The 39-compound database is a proof of concept. Scaling to $10^6$+ compounds requires engineering the trie implementation (memory management, disk-backed nodes, distributed storage). The $O(k)$ search guarantee holds regardless of $N$, but the practical implementation must handle large tries efficiently.

\item \textbf{Rotational structure.} Rotational constants are used for diatomic $\St$ computation but are not systematically included for polyatomics. A more complete treatment would incorporate rotational constants for all molecules, providing additional timescale information.

\item \textbf{Empirical parameters.} The harmonic proximity threshold $\delta = 0.05$ and the maximum ratio order $p_{\max} = 8$ are chosen empirically. A principled derivation from the bounded phase space axiom, perhaps through the resolution of the ternary partition at a given depth, would be desirable.

\item \textbf{Structural isomers.} Structural isomers with identical vibrational frequencies (if any existed) would map to the same S-entropy coordinates. In practice, different molecular structures produce distinct vibrational spectra, so this is unlikely to be a practical concern. Conformational isomers with very similar spectra may map to nearby but distinct points.

\item \textbf{Reference constants.} The normalisation constants $\omega_{\mathrm{ref,max}} = 4401$ cm$^{-1}$ (H$_2$), $\omega_{\mathrm{ref,min}} = 218$ cm$^{-1}$ (CCl$_4$), and $\Brot^{\mathrm{ref}}_{\mathrm{min}} = 0.39$ cm$^{-1}$ are determined by the extremes of the current database. If the database is extended to include molecules with higher or lower frequencies, these reference values would need to be updated, and all existing encodings would shift. A more robust approach would use physically motivated universal reference values (e.g., the Planck frequency as an upper bound), but this would reduce the effective resolution for the frequency range of interest.

\item \textbf{Degeneracy of modes.} Many molecules have degenerate vibrational modes (e.g., the $t_2$ modes of CH$_4$ are triply degenerate). The current encoding uses only distinct frequencies, not their degeneracies. Including degeneracy information would provide additional discrimination but would require extending the S-entropy definition.
\end{enumerate}

\subsection{Falsifiable Predictions}

The categorical compound database framework makes the following testable predictions:

\begin{enumerate}
\item \textbf{Spectral proximity implies ternary proximity.} Molecules with similar vibrational spectra (e.g., isotopologues, homologous series) must share ternary prefixes. The length of the shared prefix should be predictable from the spectral similarity.

\item \textbf{Chemical families must show cohesion.} For any chemically meaningful grouping (functional group families, homologous series, coordination compounds with the same ligand set), the cohesion ratio $R$ should exceed 1.0, provided the grouping reflects genuine spectral similarity.

\item \textbf{Scaling independence.} For any database of $N$ compounds, trie search at depth $k$ must complete in $O(k)$ time, independent of $N$. This is directly measurable by timing queries on databases of increasing size.

\item \textbf{Invertibility.} Given S-entropy coordinates $(\Sk, \St, \Se)$, the vibrational spectrum should be recoverable up to the resolution set by trit depth. Specifically, the number of modes and their approximate frequency distribution should be deducible from the three coordinates.

\item \textbf{Overtone enrichment.} Including overtones and combination bands in the frequency list should improve $\Se$ discrimination (more pairs to evaluate, more harmonic relationships to detect) without changing the encoding framework.

\item \textbf{Isotopologue prediction.} Isotopic substitution systematically shifts vibrational frequencies by a factor of $\sqrt{\mu/\mu'}$ where $\mu, \mu'$ are the reduced masses. The S-entropy coordinates of an isotopologue should therefore be predictable from the parent molecule's coordinates. For example, D$_2$O should have S-entropy coordinates intermediate between H$_2$O and the reference values, with the shift proportional to the square root of the mass ratio.
\end{enumerate}

Each of these predictions is falsifiable by straightforward experimental or computational tests. The framework makes no protected claims: if spectrally similar molecules fail to cluster in ternary space, or if trie search times grow with database size, the framework is refuted.

%==============================================================================
\section{Conclusion}
\label{sec:conclusion}
%==============================================================================

We have developed the categorical compound database: a molecular identification and search system in which molecules are addressed by ternary strings in a three-dimensional S-entropy coordinate space derived from vibrational spectra. The key results are:

\begin{enumerate}
\item \textbf{Mathematical foundations.} From the Bounded Phase Space Law, we derived oscillatory necessity (Theorem~\ref{thm:oscillatory}), proved that bounded molecular systems possess discrete frequency spectra (Corollary~\ref{cor:discrete}), and defined three independent S-entropy coordinates---knowledge entropy $\Sk$, temporal entropy $\St$, evolution entropy $\Se$---that characterise molecular oscillatory structure (Theorem~\ref{thm:dimensional}).

\item \textbf{Ternary encoding.} We proved that base-3 is the natural encoding for the 3D S-entropy space (Theorem~\ref{thm:ternary_natural}), established a bijection between trit strings and phase space cells (Theorem~\ref{thm:bijection}), and showed that the encoding preserves distance (Theorem~\ref{thm:distance}) and converges to continuous coordinates (Theorem~\ref{thm:continuous}).

\item \textbf{$O(k)$ search.} The ternary trie provides search in $O(k)$ operations for $k$-trit resolution, independent of database size $N$ (Theorem~\ref{thm:scaling}). For the current database, this represents a $3{,}328\times$ speedup over brute-force fingerprint comparison. At PubChem scale, the projected speedup exceeds $10^9$.

\item \textbf{Fuzzy search.} Prefix matching provides adjustable-resolution fuzzy search (Theorem~\ref{thm:fuzzy}), with each trit corresponding to one oscillation-counting observation. No similarity threshold needs to be tuned.

\item \textbf{Chemical validation.} Of six chemically defined families, five show ternary cohesion ($R > 1.0$), with chemical groupings emerging at depth 3 from S-entropy coordinates alone. All 39 compounds are uniquely resolved at depth 12.

\item \textbf{Paradigm shift.} Molecular search is reframed from \emph{computation} (evaluate a similarity function) to \emph{addressing} (traverse a phase space partition). The address is the search.
\end{enumerate}

The traditional question in molecular informatics---``how similar are these two molecules?''---is replaced by a different question: ``how many oscillation-counting observations distinguish them?'' The answer is a ternary prefix length, which simultaneously serves as the similarity metric and the search algorithm. Computation, measurement, and search collapse into a single operation: counting.

The categorical compound database is a proof of concept, but the principle it demonstrates is general. Any physical system that can be characterised by a bounded set of oscillatory modes---proteins (normal mode analysis), materials (phonon spectra), stellar atmospheres (absorption lines)---can be encoded in an analogous S-entropy coordinate space and addressed by a ternary trie. The requirement is not that the system be a small molecule, but that it be a bounded oscillatory system. This is precisely the class of systems to which the Bounded Phase Space Law applies.

The broader implication is that the right representation makes computation unnecessary. When molecular identity is encoded in the same coordinate system that governs molecular dynamics, search becomes navigation, comparison becomes co-location, and similarity becomes shared history. The categorical compound database is an instantiation of this principle: the ternary address is both the molecular identity and the search key, and the act of encoding is the act of measurement.

\subsection{Future Directions}

Several extensions of the categorical compound database are natural:

\textbf{Extended molecular coverage.} The 39-compound database should be extended to include the full NIST CCCBDB (approximately 1800 molecules with computed vibrational frequencies) and eventually to cover the spectral data available through the NIST Chemistry WebBook \citep{linstrom2023}. This would test whether the S-entropy encoding remains discriminating as the database grows, and whether chemical family cohesion is maintained at larger scales.

\textbf{Electronic spectra.} The current encoding uses only vibrational frequencies. Electronic transition energies provide an orthogonal characterisation of molecular identity. An extended S-entropy space incorporating electronic transitions would have higher dimensionality, potentially requiring a higher-base encoding or a product of independent ternary tries.

\textbf{Machine learning integration.} The S-entropy coordinates could serve as features for molecular property prediction. Unlike Morgan fingerprints, which are opaque binary vectors, the S-entropy coordinates are interpretable: $\Sk$ measures spectral complexity, $\St$ measures timescale span, and $\Se$ measures harmonic connectivity. These physically meaningful features may improve the interpretability of machine learning models for molecular properties.

\textbf{Reaction pathway search.} If reactants and products are both encoded in S-entropy space, a chemical reaction corresponds to a displacement vector in $\Sspace$. Reaction types that conserve certain S-entropy relationships (e.g., reactions that preserve harmonic connectivity) would correspond to constrained paths in S-entropy space. This could enable a ``categorical retrosynthesis'' in which the search for synthetic routes is guided by S-entropy topology rather than by structural transformation rules.

\textbf{Experimental validation.} The S-entropy encoding can be validated directly against experimental IR and Raman spectra. Given a measured spectrum, the encoding protocol produces an address, and the trie returns candidate identifications. This could be implemented as a spectral database search tool that bypasses the traditional spectrum-to-structure-to-fingerprint pipeline entirely.

\section*{Data Availability}

All vibrational frequencies used in this work are from the NIST Computational Chemistry Comparison and Benchmark Database \citep{johnson2022}. The complete Python implementation of the categorical compound database, including the encoding protocol, trie construction, fuzzy search, and validation analysis, is available at \url{https://github.com/kundai-sachikonye/borgia}.

\section*{Acknowledgements}

The author thanks the NIST CCCBDB for providing freely accessible vibrational frequency data.

\bibliographystyle{plainnat}
\bibliography{references}

\end{document}